\chapter{环面几何}\label{chap:toric}
\SourcePage{54}

本章将研究复数域上的环面簇理论。我们假定读者对用锥与扇描述环面簇的经典方法已有一定了解，参见 \SourceCite{Danilov,Fulton3,Oda}；不过，我们也会介绍较近出现的、利用多胞形与齐次坐标构造环面簇的方法。我们还将讨论环面几何中与镜像对称关系最密切的内容，包括 Kähler 锥、辛约化、GKZ 分解、Fano 环面簇与自反多胞形，以及自同构群。本章最后给出一些例子。

\section{锥与扇}\label{sec:cones-fans}

首先概述环面簇，主要目的是确定记号和术语。有关证明可参见 \SourceCite{Danilov,Fulton3,Oda}。

设 $M\simeq\mathbb Z^n$ 是秩为 $n$ 的自由 Abel 群，$N=\operatorname{Hom}(M,\mathbb Z)$ 是它的对偶群。$m\in M$ 与 $v\in N$ 的配对记为 $\langle m,v\rangle\in\mathbb Z$。$N_{\mathbb R}=N\otimes\mathbb R$ 中的\emph{有理多面锥} $\sigma$ 是形如
\[
\sigma=\left\{\sum_{i=1}^{s}\lambda_i u_i:\lambda_i\geq0\right\}
\]
的子集，其中 $u_1,\ldots,u_s\in N$。若 $\sigma\cap(-\sigma)=\{0\}$，则称 $\sigma$\emph{强凸}。$\sigma$ 的维数是它在 $N_{\mathbb R}$ 中张成的子空间的维数。每个锥 $\sigma$ 都有\emph{对偶锥}
\[
\check\sigma=\{m\in M_{\mathbb R}:\langle m,v\rangle\geq0\text{ 对所有 }v\in\sigma\}.
\]
标准理论表明，$\check\sigma$ 也是有理多面锥。给定 $m\in M\cap\check\sigma$，子集
\[
\tau=\{v\in\sigma:\langle m,v\rangle=0\}\subset\sigma
\]
称为 $\sigma$ 的一个\emph{面}。每个面仍是有理多面锥；面的面仍是原锥的面。

$N_{\mathbb R}$ 中的\emph{扇} $\Sigma$ 是一个由有限多个强凸有理多面锥组成的集合，满足：
\begin{itemize}
\item 若 $\sigma\in\Sigma$，则 $\sigma$ 的每个面也属于 $\Sigma$；
\item 若 $\sigma,\tau\in\Sigma$，则 $\sigma\cap\tau$ 是二者各自的一个面。
\end{itemize}
集合 $|\Sigma|=\bigcup_{\sigma\in\Sigma}\sigma\subset N_{\mathbb R}$ 称为 $\Sigma$ 的\emph{支集}。对每个 $d$，用 $\Sigma(d)$ 表示 $\Sigma$ 中所有 $d$ 维锥的集合。

一维锥有时称为扇的\emph{一维骨架}，尤为重要。我们约定用 $\rho$ 表示 $\Sigma(1)$ 中的元素。对每个 $\rho$，令 $v_\rho$ 为半群 $\rho\cap N$ 的唯一生成元。借助这些生成元，$\Sigma$ 中的锥 $\sigma$ 可写为
\[
\sigma=\left\{\sum_{\rho\subset\sigma}\lambda_\rho v_\rho:\lambda_\rho\geq0\right\}.
\]

\SourcePage{55}
属于 $\sigma$ 的各 $v_\rho$ 称为 $\sigma$ 的生成元。按照惯常的记号简化，我们常把 $\rho$ 与 $v_\rho$ 等同起来，于是也写 $v_\rho\in\Sigma(1)$。若 $r=|\Sigma(1)|$ 是一维锥的数目，有时还把这些 $v_\rho$ 记为 $v_1,\ldots,v_r$。\index{扇|hyperpage}\index{有理多面锥|hyperpage}\index{环面簇|hyperpage}

下面回顾环面簇的经典定义。给定扇 $\Sigma$，每个锥 $\sigma\in\Sigma$ 都给出仿射环面簇
\[
X_\sigma=\operatorname{Spec}\bigl(\mathbb C[M\cap\check\sigma]\bigr).
\]
这里，$\mathbb C[M\cap\check\sigma]$ 是这样一个 $\mathbb C$ 代数：对每个 $m\in M\cap\check\sigma$ 有生成元 $\chi^m$，并满足关系 $\chi^m\chi^{m'}=\chi^{m+m'}$。沿着 $X_{\sigma\cap\tau}$ 将各 $X_\sigma$ 与 $X_\tau$ 粘合，便得到环面簇 $X_\Sigma$。

若 $\Sigma$ 是 $N_{\mathbb R}\simeq\mathbb R^n$ 中的扇，则 $X_\Sigma$ 是维数为 $n$ 的 Cohen--Macaulay 代数簇。由于平凡锥 $\{0\}$ 对应的仿射环面簇是代数环面 $T_N=N\otimes\mathbb C^*=\operatorname{Spec}(\mathbb C[M])$，故 $T_N$ 是 $X_\Sigma$ 的仿射开子集。环面在自身上的作用可延拓为 $T_N$ 在 $X_\Sigma$ 上的代数作用，“环面簇”由此得名。在不会混淆时，我们简写 $X$，不写 $X_\Sigma$。

环面在 $X$ 上的作用只有有限多个轨道，轨道闭包与 $\Sigma$ 中的锥之间存在反转包含关系的一一对应。特别地，每个 $\rho\in\Sigma(1)$ 对应一个不可约的 $T_N$ 不变除子 $D_\rho\subset X$。此外，$m\in M$ 给出一个特征标 $\chi^m:T_N\to\mathbb C^*$。把 $\chi^m$ 看作 $X$ 上的有理函数，就有
\begin{equation}\label{eq:toric-principal-divisor}
\operatorname{div}(\chi^m)=\sum_\rho\langle m,v_\rho\rangle D_\rho .
\end{equation}

扇 $\Sigma$ 的性质对环面簇 $X$ 的几何性质有深刻影响。例如：
\begin{itemize}
\item $X$ 是完备的，即紧的，当且仅当 $|\Sigma|=N_{\mathbb R}$；这样的扇称为\emph{完备扇}。
\item $X$ 光滑，即无奇点，当且仅当 $\Sigma$ 中每个锥的生成元都是 $N$ 的某组 $\mathbb Z$ 基的一部分；这样的扇称为\emph{光滑扇}。
\item $X$ 是 轨形，即 $V$ 流形，当且仅当 $\Sigma$ 中每个锥的生成元在 $\mathbb R$ 上线性无关；此时称 $\Sigma$ 和 $X$ 为\emph{单纯}的。对环面簇 $X$，“$X$ 单纯 $\Rightarrow$ $X$ 是轨形”是初等的；反之，“$X$ 是轨形 $\Rightarrow$ $X$ 单纯”则深刻得多，其证明见 M.~Brion, \emph{Rational smoothness and fixed points of torus actions}, Transformation Groups 4 (1999), 127--156。
\end{itemize}
一般的环面簇只有较温和的奇点，因为它是 Cohen--Macaulay 的；不过，我们有时仍会要求 $X$ 单纯。\hyperref[app:singular-varieties]{附录 A} 将说明，\IndexLocation{indexloc:1375:55}单纯环面簇在许多方面具有流形的性质。

由 Weil 除子模去线性等价所得的 Chow 群 $A_{n-1}(X)$，可直接通过扇计算。我们始终假定一维锥张成 $N_{\mathbb R}$。由式~\eqref{eq:toric-principal-divisor} 得到正合列
\begin{equation}\label{eq:toric-chow-sequence}
0\longrightarrow M\longrightarrow\mathbb Z^{\Sigma(1)}
\longrightarrow A_{n-1}(X)\longrightarrow0 .
\end{equation}
其中 $m\in M$ 映到 $(\langle m,v_\rho\rangle)\in\mathbb Z^{\Sigma(1)}$，而 $(a_\rho)\in\mathbb Z^{\Sigma(1)}$ 映到除子 $\sum_\rho a_\rho D_\rho$ 的类。因此 $A_{n-1}(X)$ 的秩为 $r-n$。

$A_{n-1}(X)$ 中包含 Picard 群 $\operatorname{Pic}(X)$，它由 Cartier 除子模去线性等价构成。Weil 除子 $D=\sum_\rho a_\rho D_\rho$ 是 Cartier 除子，当且仅当对每个 $\sigma\in\Sigma$，都有 $m_\sigma\in M$，使得只要 $\rho\subset\sigma$ 就有 $\langle m_\sigma,v_\rho\rangle=-a_\rho$。函数 $\phi_D:|\Sigma|\to\mathbb R$，在 $v\in\sigma$ 时定义为 $\phi_D(v)=\langle m_\sigma,v\rangle$，称为 $D$ 的\emph{\IndexLocation{indexloc:1314:55}支撑函数}。

若 $X$ 光滑，则 $\operatorname{Pic}(X)=A_{n-1}(X)$。在单纯情形，Picard 群是 Chow 群的有限指数子群。$X$ 完备时 Picard 群总是无挠的；而 $A_{n-1}(X)$ 即使在 $X$ 单纯时也可能有挠。

\SourcePage{56}
若 $X$ 完备且除子 $D=\sum_\rho a_\rho D_\rho$ 是 Cartier 除子，可按以下条件判定 $D$ 是否丰沛，或是否由整体截面生成：
\begin{itemize}
\item $D$ 由整体截面生成，当且仅当对任意 $\rho\not\subset\sigma$，均有 $\langle m_\sigma,v_\rho\rangle\geq-a_\rho$。
\item $D$ 丰沛，当且仅当对任意 $n$ 维锥 $\sigma$ 及 $\rho\not\subset\sigma$，均有 $\langle m_\sigma,v_\rho\rangle>-a_\rho$。
\end{itemize}
第一种情形称 $\phi_D$\emph{上凸}，第二种情形称它\emph{严格上凸}。下一节将用某些多胞形解释这些条件。

\section{多胞形与齐次坐标}\label{sec:polytopes-coordinates}

上一节给出的环面簇经典定义，需要将仿射环面簇粘合起来。近来又发现了其他构造方法。本节研究其中两种：一种来自多胞形，另一种推广了 $\mathbb P^n$ 的通常齐次坐标。涉及辛约化的第三种构造将在\hyperref[subsec:toric-symplectic]{第 3.3.3 小节}讨论。

\subsection{多胞形}\label{subsec:toric-polytopes}

先确定术语。$M_{\mathbb R}$ 中的\emph{多胞形} $\Delta$ 是有限个点的凸包。$\Delta$ 的维数，是所有差 $m_1-m_2$（$m_1,m_2\in\Delta$）所张成子空间的维数。若 $\Delta$ 的顶点都属于 $M$，即 $\Delta$ 是 $M$ 中有限子集的凸包，则称其为\emph{整多胞形}。\emph{刻面}是余维数为一的面。

给定 $M_{\mathbb R}$ 中的多胞形 $\Delta_1,\ldots,\Delta_k$，可以用几种方法构造新多胞形。例如，它们并集的凸包记为
\[
\operatorname{Conv}(\Delta_1,\ldots,\Delta_k).
\]
还可定义\emph{Minkowski 和}：
\[
\Delta_1+\cdots+\Delta_k
=\{m_1+\cdots+m_k:m_i\in\Delta_i\}.
\]
若 $k\Delta=\Delta+\cdots+\Delta$（共 $k$ 项），则由凸性可知 $k\Delta=\{km:m\in\Delta\}$。\hyperref[sec:toric-complete-intersections]{第 4.3 节}将讨论的镜像对称构造，关键地用到了凸包与 Minkowski 和。

研究环面簇时，多胞形自然出现。例如，若 $X$ 完备，$D=\sum_\rho a_\rho D_\rho$ 是 Cartier 除子，则
\begin{equation}\label{eq:divisor-polytope}
\begin{aligned}
\Delta_D
&=\{m\in M_{\mathbb R}:\langle m,v\rangle\geq\phi_D(v)
\text{ 对所有 }v\in N_{\mathbb R}\}\\
&=\{m\in M_{\mathbb R}:\langle m,v_\rho\rangle\geq-a_\rho
\text{ 对所有 }\rho\}
\end{aligned}
\end{equation}
是一个多胞形。\footnote{即使 $D$ 不是 Cartier 除子，也可用第二个等式定义 $\Delta_D$。}
此外，$\Delta_{kD}=k\Delta_D$，$\Delta_{D+\operatorname{div}(\chi^m)}=\Delta_D-m$，并且：
\begin{itemize}
\item $D$ 由整体截面生成，当且仅当 $\Delta_D$ 是集合 $\{m_\sigma:\sigma\in\Sigma(n)\}$ 的凸包；
\item $D$ 丰沛，当且仅当对 $\Sigma(n)$ 中不同的 $\sigma,\tau$ 有 $m_\sigma\ne m_\tau$，且 $\Delta_D$ 是以 $\{m_\sigma:\sigma\in\Sigma(n)\}$ 为顶点集合的 $n$ 维多胞形。
\end{itemize}
在这两种情形，$\Delta_D$ 都是整多胞形。此外，当 $D$ 丰沛时，$\Delta_D$ 的非空面与 $\Sigma$ 中的锥一一对应。特别地，$\Delta_D$ 的刻面对应于 $\Sigma(1)$ 中的元素。这些断言的证明可参见 \SourceCite{Danilov,Fulton3,Oda}。\index{多胞形|hyperpage}\index{Minkowski和@Minkowski 和|hyperpage}

\SourcePage{57}
另一个有用的事实是，$\Delta_D\cap M$ 与 $H^0(X,\mathcal O(D))$ 中的 $T_N$ 特征向量自然对应。换言之，有一个 $T_N$ 等变同构
\begin{equation}\label{eq:toric-sections}
H^0(X,\mathcal O(D))\simeq
\bigoplus_{m\in\Delta_D\cap M}\mathbb C\chi^m .
\end{equation}
要理解这个等式，可将 $\mathcal O(D)$ 的截面看作 $X$ 上满足 $\operatorname{div}(f)+D\geq0$ 的有理函数 $f$。由式~\eqref{eq:toric-principal-divisor}，$\operatorname{div}(\chi^m)+D\geq0$ 等价于对所有 $\rho$ 均有 $\langle m,v_\rho\rangle\geq-a_\rho$。上述同构由此很容易推出，详情参见 \SourceCite[第 3.4 节]{Fulton3}。

\subsection{由多胞形构造环面簇}\label{subsec:toric-from-polytopes}

下面说明如何从一个 $n$ 维整多胞形出发构造环面簇。我们采用 \SourceCite{Batyrev4} 的方法。给定 $\Delta$，考虑 $m\in k\Delta\cap M$ 时的“单项式” $t_0^k\chi^m$。它们按规则
$t_0^k\chi^m\cdot t_0^\ell\chi^{m'}
=t_0^{k+\ell}\chi^{m+m'}$
相乘；因为 $m\in k\Delta\cap M$、$m'\in\ell\Delta\cap M$ 意味着 $m+m'\in(k+\ell)\Delta\cap M$。由这些 $t_0^k\chi^m$ 生成的 $\mathbb C$ 代数记为 $S_\Delta$，并规定
\[
\deg(t_0^k\chi^m)=k
\]
来赋予它分次。我们称 $S_\Delta$ 为 $\Delta$ 的\emph{多胞形环}。令
\[
\mathbb P_\Delta=\operatorname{Proj}(S_\Delta)
\]
为相应的射影簇。在语境清楚时，用 $\mathbb P$ 简记 $\mathbb P_\Delta$。

要把 $\mathbb P=\mathbb P_\Delta$ 描述为环面簇，须给出它在 $N_{\mathbb R}$ 中的扇。给定非空面 $F\subset\Delta$，考虑锥
\[
\check\sigma_F
=\{\lambda(m-m'):m\in\Delta,\ m'\in F,\ \lambda\geq0\}
\subset M_{\mathbb R}.
\]
它的对偶是 $N_{\mathbb R}$ 中的锥 $\sigma_F$。所有这些锥组成扇
$\{\sigma_F:F\text{ 是 }\Delta\text{ 的非空面}\}$，称为 $\Delta$ 的\emph{法扇}。这是 $N_{\mathbb R}$ 中的完备扇，而且可以证明 $\mathbb P$ 正是由该法扇确定的环面簇 \SourceCite{Batyrev4}。\index{法扇|hyperpage}

当原点是 $\Delta$ 的内点时，法扇特别容易形象化。先定义\emph{\IndexLocation{indexloc:1077:57}极多胞形}
\[
\Delta^\circ
=\{v\in N_{\mathbb R}:\langle m,v\rangle\geq-1
\text{ 对所有 }m\in\Delta\}\subset N_{\mathbb R}.
\]
$\Delta^\circ$ 也以原点为内点，且 $(\Delta^\circ)^\circ=\Delta$。此外，$\Delta$ 的每个 $i$ 维面 $F$ 对应于 $\Delta^\circ$ 的一个 $(n-1-i)$ 维面 $F^\circ$；其特征是，当 $m'\in F$、$v\in F^\circ$ 时有 $\langle m',v\rangle=-1$。利用 $\Delta^\circ$，可如下描述 $\mathbb P$。

\begin{lemma}\label{lem:normal-polar}
在 $N_{\mathbb R}$ 中，由 $\Delta^\circ$ 的各真面上的锥组成的扇，就是 $\Delta$ 的法扇，因此给出环面簇 $\mathbb P=\mathbb P_\Delta$。
\end{lemma}
\begin{proof}
只须证明：对 $\Delta$ 的每个真面 $F$，
$\sigma_F=\{\lambda v:v\in F^\circ,\ \lambda\geq0\}$。
首先，$v\in\sigma_F$ 当且仅当对所有 $m\in\Delta$ 和 $m'\in F$，都有 $\langle m,v\rangle\geq\langle m',v\rangle$。还要注意，若 $v\ne0$，则 $\langle m',v\rangle<0$；否则就会对所有 $m\in\Delta$ 有 $\langle m,v\rangle\geq0$，这与原点是 $\Delta$ 的内点矛盾。

若 $v\in F^\circ$，则 $\langle m,v\rangle\geq\langle m',v\rangle=-1$，所以容易得到 $\lambda v\in\sigma_F$。反过来，设对所有 $m\in\Delta$、$m'\in F$ 均有 $\langle m,v\rangle\geq\langle m',v\rangle$。可设 $v\ne0$。固定 $F$ 相对内部的一个 $m'$；由上面的观察，$\langle m',v\rangle<0$。因此可写 $v=\lambda v'$，其中 $\lambda>0$，并且对所有 $m\in\Delta$ 有 $\langle m,v'\rangle\geq\langle m',v'\rangle=-1$。于是 $v'\in\Delta^\circ$；又因 $m'$ 在 $F$ 的相对内部，$\langle m',v'\rangle=-1$ 进一步表明 $v'\in F^\circ$。
\end{proof}

\SourcePage{58}
由引理~\ref{lem:normal-polar} 可知，在 $\mathbb P_\Delta$ 中，$T_N$ 轨道闭包与 $\Delta^\circ$ 的真面一一对应；再由极对偶性，它们也与 $\Delta$ 的非空面 $F$ 一一对应。实际上，从多胞形环出发也能直接看出这种对应。

设 $F\subset\Delta$ 是一个面。多胞形环之间有典范满射 $\pi_F:S_\Delta\to S_F$，定义为
\[
\pi_F(t_0^k\chi^m)=
\begin{cases}
t_0^k\chi^m,&m\in kF,\\
0,&m\in k\Delta-kF.
\end{cases}
\]
这个满射诱导环面簇的自然嵌入 $\mathbb P_F\hookrightarrow\mathbb P_\Delta$。若 $S_\Delta$ 由 $\Delta\cap M$ 生成，这个嵌入可以描述得特别简单：可把 $\mathbb P_\Delta$ 看成嵌入一个射影空间，其齐次坐标与 $\Delta\cap M$ 中的点一一对应。此时，只需把对应于 $(\Delta-F)\cap M$ 中各点的坐标全部置零，便得到 $\mathbb P_F$。

若从一个以原点为内点的整多胞形 $\Delta$ 出发，其极多胞形 $\Delta^\circ$ 未必是整的。\hyperref[sec:fano-reflexive]{第 3.5 节}将研究 Batyrev 的一类特殊多胞形，即极多胞形也是整的\emph{自反多胞形}。这类多胞形在许多镜像对称构造中起重要作用。\index{自反多胞形|hyperpage}

再回到任意 $n$ 维整多胞形 $\Delta$。我们注意到，环面簇 $\mathbb P=\mathbb P_\Delta$ 自然带有一个丰沛线丛：$\mathbb P=\operatorname{Proj}(S_\Delta)$ 带有丰沛的 $\mathcal O_{\mathbb P}(1)$。从环面簇的观点，这个线丛可描述如下。定义函数 $\phi_\Delta:N_{\mathbb R}\to\mathbb R$：
\[
\phi_\Delta(v)=\min_{m\in\Delta}\langle m,v\rangle .
\]
可以证明，$\phi_\Delta$ 是 $\mathbb P$ 上某个除子 $D_\Delta$ 的支撑函数。而且，$\phi_\Delta$ 严格上凸，$\mathcal O_{\mathbb P}(D_\Delta)\simeq\mathcal O_{\mathbb P}(1)$，与 $D_\Delta$ 对应的多胞形 $\Delta_{D_\Delta}$ 恰好是 $\Delta$，参见 \SourceCite{Batyrev4,Oda}。所以，由多胞形构造出的环面簇，是附带选定丰沛除子的射影簇。

若选定 $M$ 的一组基，就能更具体地观察 $\mathbb P_\Delta$ 及其射影嵌入。这等于为环面 $T_N$ 选取坐标 $t_1,\ldots,t_n$。当 $m\in M$ 写为 $m=(a_1,\ldots,a_n)$ 时，有 $\chi^m=\prod_{i=1}^{n}t_i^{a_i}$，因此可以写 $t^m$ 来代替 $\chi^m$。给定多胞形 $\Delta$，选择 $k$ 使 $kD_\Delta$ 十分丰沛。格点 $k\Delta\cap M=\{m_1,\ldots,m_s\}$ 给出单项式 $t^{m_i}$，从而定义映射 $T_N\to\mathbb P^{s-1}$，把 $t\in T_N$ 送到 $(t^{m_1},\ldots,t^{m_s})$。$\mathbb P_\Delta$ 就是此映射的像在 $\mathbb P^{s-1}$ 中的闭包。

另外，对任意 $k\geq0$，有 Laurent 多项式的向量空间
\begin{equation}\label{eq:laurent-space}
L(k\Delta)=\left\{f:f=\sum_{m\in k\Delta\cap M}a_mt^m,\ a_m\in\mathbb C\right\}.
\end{equation}
每个 $f\in L(k\Delta)$ 都定义了 $T_N$ 中由 $f=0$ 给出的仿射超曲面 $Z_f$；其在 $\mathbb P_\Delta$ 中的闭包 $\overline Z_f$ 是相应的除子。这样，把 $M$ 中的元素视为 Laurent 单项式的指数向量，便可用坐标研究 $\mathbb P_\Delta$。有时人们从 Laurent 多项式 $f$ 出发，将 $\Delta$ 定义为其指数向量的凸包，这就是 $f$ 的\emph{Newton 多胞形}。我们使用以下记号：
\begin{equation}\label{eq:lattice-points}
\begin{aligned}
l(k\Delta)&=|k\Delta\cap M|=\dim L(k\Delta),\\
l^*(k\Delta)&=|\{m\in k\Delta\cap M:m\text{ 不在 }k\Delta\text{ 的任何刻面上}\}|.
\end{aligned}
\end{equation}

\SourcePage{59}
数 $l(k\Delta)$ 和 $l^*(k\Delta)$ 具有许多有趣的性质，并在镜像对称的多种情境中出现。

最后对前面定义的多胞形环 $S_\Delta$ 作几点说明。首先，对 $f\in L(k\Delta)$ 而言，多项式 $t_0^kf$ 恰好是 $S_\Delta$ 中次数为 $k$ 的元素。这给出同构
\begin{equation}\label{eq:polytope-ring-sections}
(S_\Delta)_k\simeq L(k\Delta)\simeq
H^0(\mathbb P,\mathcal O_{\mathbb P}(kD_\Delta)),
\end{equation}
其中第二个同构由式~\eqref{eq:toric-sections} 得到。

还可将 $S_\Delta$ 解释为仿射坐标环。令 $\check\sigma_\Delta\subset\mathbb R\times M_{\mathbb R}$ 为 $\{1\}\times\Delta$ 上的锥。注意 $(k,m)\in\check\sigma_\Delta\cap(\mathbb Z\times M)$ 当且仅当 $m\in k\Delta$。因此，若 $\sigma_\Delta\subset\mathbb R\times N_{\mathbb R}$ 是 $\check\sigma_\Delta$ 的对偶，则 $S_\Delta$ 正是由 $\sigma_\Delta$ 给出的仿射环面簇的坐标环。由 \SourceCite{Danilov} 可知 $S_\Delta$ 是 Cohen--Macaulay 的。

\subsection{齐次坐标环}\label{subsec:homogeneous-ring}

紧接在式~\eqref{eq:laurent-space} 之前的讨论中，我们在环面簇的环面上引入了坐标 $t_1,\ldots,t_n$。这些坐标非常有用；但若能像射影空间上一样拥有整体坐标，就更方便了。为此，我们定义环面簇的齐次坐标环。

若 $X=X_\Sigma$ 由 $N_{\mathbb R}$ 中的扇 $\Sigma$ 给出，就为每个 $\rho\in\Sigma(1)$ 引入变量 $x_\rho$，并考虑多项式环
\[
S=\mathbb C[x_\rho:\rho\in\Sigma(1)].
\]
$S$ 中的单项式写作 $x^D=\prod_\rho x_\rho^{a_\rho}$，其中 $D=\sum_\rho a_\rho D_\rho$ 是 $X$ 上有效的环面不变除子；这里用到了对应 $\rho\leftrightarrow D_\rho$。规定 $x^D$ 的次数为
\[
\deg(x^D)=[D]\in A_{n-1}(X).
\]
于是 $S$ 按 Chow 群 $A_{n-1}(X)$ 分次。给定除子类 $\alpha\in A_{n-1}(X)$，以 $S_\alpha$ 表示 $S$ 中次数为 $\alpha$ 的分次部分。我们也常把变量写成 $x_1,\ldots,x_r$，其中 $x_i$ 对应锥生成元 $v_i$，$r=|\Sigma(1)|$；这时 $S=\mathbb C[x_1,\ldots,x_r]$。

\begin{definition}\label{def:toric-homogeneous-ring}
环 $S$ 连同上面定义的分次，称为 $X$ 的\emph{齐次坐标环}。
\end{definition}

\hyperref[sec:toric-examples]{第 3.7 节}将看到，$\mathbb P^n$ 的齐次坐标环就是带通常分次的 $\mathbb C[x_0,\ldots,x_n]$。对 $\mathbb P^n\times\mathbb P^m$，坐标环则是按双齐次多项式的相应次数分次的环。在与镜像对称有关的物理理论中，变量 $x_\rho\in S$ 对应某些可观测量，尽管单个可观测量本身并没有物理意义。所使用的是由它们组成的某些单项式，再模去 $S$ 的一个理想；这就是\emph{手征环}。

注意，分次环 $S$ 只依赖于 $\Sigma$ 的一维骨架。因此，可能存在许多具有同一一维骨架的扇，从而有许多具有同一坐标环的环面簇。\hyperref[sec:gkz-decomposition]{第~\ref*{sec:gkz-decomposition} 节}将进一步阐明这一点。

为在 $S$ 中体现整个扇 $\Sigma$，定义理想 $B(\Sigma)$。对 $\sigma\in\Sigma$，令
\[
\widehat x_\sigma=\prod_{\rho\not\subset\sigma}x_\rho
\]
为不属于 $\sigma$ 的变量之积，再定义理想
\[
B(\Sigma)=(\widehat x_\sigma:\sigma\in\Sigma)\subset S.
\]

\SourcePage{60}
对 $\mathbb P^n$，这个理想就是“无关”理想 $\langle x_0,\ldots,x_n\rangle$。一般而言，$B(\Sigma)$ 唯一地确定 $\Sigma$。

$S$ 的各分次部分还有一个很好的上同调解释。若除子 $D=\sum_\rho a_\rho D_\rho$ 不是 Cartier 除子，它就不能确定线丛，但仍给出一个自反层 $\mathcal O_X(D)$。这是 $X$ 上有理函数层的子层：它在开集 $U\subset X$ 上的截面，是在 $U$ 上极点阶数不超过 $D$ 所规定阶数的有理函数。关于 Weil 除子所确定的层，详见 \SourceCite[第 1 节附录]{Reid1}。如式~\eqref{eq:polytope-ring-sections} 所示，有同构
\[
H^0(X,\mathcal O_X(D))\simeq L(\Delta_D),
\]
这里 $\Delta_D$ 由式~\eqref{eq:divisor-polytope} 的第二行定义。由此容易看出，映射
$t^m\mapsto\prod_\rho x_\rho^{\langle m,v_\rho\rangle+a_\rho}$
诱导同构
\begin{equation}\label{eq:toric-graded-sections}
H^0(X,\mathcal O_X(D))\simeq S_\alpha,
\end{equation}
其中 $\alpha=[D]\in A_{n-1}(X)$。

在射影情形，可将它与式~\eqref{eq:polytope-ring-sections} 的同构结合。于是，若 $D$ 丰沛且 $\alpha=[D]$，就得到环同构
\[
S_{\Delta_D}\simeq\bigoplus_{k=0}^{\infty}S_{k\alpha};
\]
这里 $S_{\Delta_D}$ 是\hyperref[subsec:toric-from-polytopes]{第 3.2.1 小节}的多胞形环。因此，$S$ 包含了来自 $X$ 上所有可能丰沛除子 $D$ 的多胞形环 $S_{\Delta_D}$。

还可以用齐次坐标描述 $X$ 上的微分形式。\hyperref[chap:hodge-yukawa]{第 5 章}计算 $X$ 中丰沛超曲面的 Yukawa 耦合时，将采用这一描述。

\subsection{由齐次坐标构造环面簇}\label{subsec:toric-from-homogeneous}

下面说明怎样把射影空间的通常构造
$\mathbb P^n=(\mathbb C^{n+1}-\{0\})/\mathbb C^*$
推广到任意环面簇。如\hyperref[sec:cones-fans]{第 3.1 节}一样，假定 $\Sigma$ 的一维锥张成 $N_{\mathbb R}$；若 $\Sigma$ 完备，这一条件自动成立。定义群
\[
G=\operatorname{Hom}_{\mathbb Z}(A_{n-1}(X),\mathbb C^*).
\]
对正合列~\eqref{eq:toric-chow-sequence} 施加 $\operatorname{Hom}_{\mathbb Z}(-,\mathbb C^*)$，得到正合列
\begin{equation}\label{eq:toric-g-sequence}
1\longrightarrow G\longrightarrow(\mathbb C^*)^{\Sigma(1)}
\longrightarrow T_N\longrightarrow1.
\end{equation}
这里用到了
$\operatorname{Hom}_{\mathbb Z}(M,\mathbb C^*)
=N\otimes\mathbb C^*=T_N$。
所以，$G$ 同构于 $(\mathbb C^*)^{r-n}$ 与一个有限 Abel 群的乘积；当且仅当 $A_{n-1}(X)$ 有挠时，才出现这个有限群。于是，$G\to(\mathbb C^*)^{\Sigma(1)}$ 给出 $G$ 在 $\mathbb C^{\Sigma(1)}=\operatorname{Spec}(S)$ 上的作用。若 $a=(a_\rho)\in\mathbb C^{\Sigma(1)}$、$g\in G$，则
\[
g\cdot a=(g([D_\rho])a_\rho).
\]
按 $G$ 的定义，$g:A_{n-1}(X)\to\mathbb C^*$。另外，理想 $B(\Sigma)\subset S$ 在 $\mathbb C^{\Sigma(1)}$ 中定义簇 $Z(\Sigma)$。这些数据按下述方式确定 $X$。

\begin{theorem}\label{thm:toric-quotient}
若 $\Sigma$ 的一维锥张成 $N_{\mathbb R}$，则：
\begin{enumerate}[label=(\roman*)]
\item $X$ 是 $\mathbb C^{\Sigma(1)}-Z(\Sigma)$ 关于 $G$ 的范畴商；
\item $X$ 是 $\mathbb C^{\Sigma(1)}-Z(\Sigma)$ 关于 $G$ 的几何商，当且仅当 $X$ 单纯。
\end{enumerate}
\end{theorem}

\SourcePage{61}
在单纯情形，这个结果最早由 \SourceCite{Audin} 证明；不过，这一构造的想法曾多次独立出现。有关证明和参考文献清单，参见 \SourceCite{Cox}。

定理~\ref{thm:toric-quotient} 给出的 $X$ 的描述，也使\hyperref[sec:cones-fans]{第 3.1 节}提到的轨道闭包与 $\Sigma$ 中锥之间反转包含关系的对应变得简单：对每个 $\sigma\in\Sigma$，考虑线性子空间
\[
\{x\in\mathbb C^{\Sigma(1)}:x_\rho=0
\text{ 对所有 }\rho\in\sigma\}.
\]
它在 $X_\Sigma$ 中下降为一个子簇，容易验证这就是与 $\sigma$ 对应的轨道闭包。特别地，与边 $v_i$ 对应的环面除子 $D_i$，正是方程 $x_i=0$ 所定义的超曲面。

定理~\ref{thm:toric-quotient} 中的簇 $Z(\Sigma)$ 有简单结构。若子集 $\mathcal S\subset\Sigma(1)$ 不是 $\Sigma$ 中某个锥的全体一维锥，但 $\mathcal S$ 的每个真子集都是某个锥的全体一维锥，就称 $\mathcal S$ 为\emph{本原集合}。这时
\begin{equation}\label{eq:toric-primitive-z}
Z(\Sigma)=\bigcup_{\mathcal S}V(\mathcal S),
\end{equation}
其中并集遍历所有本原集合 $\mathcal S$，$V(\mathcal S)$ 是对 $\rho\in\mathcal S$ 令 $a_\rho=0$ 所定义的子空间。证明见 \SourceCite{BC}。

若 $X$ 单纯，定理~\ref{thm:toric-quotient} 表明我们的确有了齐次坐标：$X$ 中一点的坐标可取为
$a=(a_\rho)\in\mathbb C^{\Sigma(1)}-Z(\Sigma)$；另一个
$b\in\mathbb C^{\Sigma(1)}-Z(\Sigma)$ 表示 $X$ 中同一点，当且仅当存在 $g\in G$，使 $b=g\cdot a$。特别地，若 $f\in S$ 是齐次的，即 $f\in S_\alpha$ 对某个 $\alpha$ 成立，那么 $f=0$ 就在 $X$ 中定义一个超曲面。更一般地，$S$ 的任意齐次理想 $I$ 都定义 $X$ 中的一个子簇 $V(I)$。环面簇的\IndexLocation{indexloc:1370:61}理想与簇的对应关系见 \SourceCite{Cox}。

在非单纯情形，$\mathbb C^{\Sigma(1)}-Z(\Sigma)$ 与 $X$ 的关系不那么直接，但 $S$ 的齐次理想仍能确定 $X$ 的子簇。还有一种值得注意的现象，即使在单纯情形也会出现：齐次方程 $f=0$ 给出的超曲面 $Y\subset X$ 未必是 Cartier 除子。这意味着，虽然 $Y$ 可由整体方程 $f=0$ 定义，在 $X$ 上的局部却未必能由单个方程定义。\index{齐次坐标环|hyperpage}\index{本原集合|hyperpage}

\section{Kähler 锥与辛几何}\label{sec:kahler-symplectic}

本节描述\IndexLocation{indexloc:1375:61}单纯环面簇的 Kähler 锥，并解释如何利用辛几何构造环面簇。

\subsection{环面簇的 Kähler 锥}\label{subsec:toric-kahler-cone}

对于光滑射影环面簇 $X$，有 $A_{n-1}(X)\simeq H^2(X,\mathbb Z)$，从而 $H^2(X,\mathbb R)=H^{1,1}(X,\mathbb R)$。更一般地，若 $X$ 单纯且射影，则 $A_{n-1}(X)\otimes\mathbb R\simeq H^2(X,\mathbb R)$，所以仍有 $H^2(X,\mathbb R)=H^{1,1}(X,\mathbb R)$。回顾\hyperref[chap:introduction]{第 1 章}，$H^{1,1}(X,\mathbb R)$ 中的 Kähler 锥由 $X$ 上所有 Kähler 类组成；轨形 情形下 Kähler 的定义见\hyperref[app:singular-varieties]{附录 A}。在上述同构下，我们得到 $A_{n-1}(X)\otimes\mathbb R$ 中的一个锥，也称为 Kähler 锥。

令 $A^+_{n-1}(X)\otimes\mathbb R$ 为由除子类 $[D_i]$（$v_i\leftrightarrow D_i$，$i=1,\ldots,r$）生成的锥。Kähler 锥包含在其中，描述如下。设
$a=\sum_{i=1}^{r}a_i[D_i]\in A^+_{n-1}(X)\otimes\mathbb R$。
因为 $X$ 单纯，对每个 $\sigma\in\Sigma$ 都可找到 $m_\sigma\in M_{\mathbb R}$，使得在 $v_i\in\sigma$ 时 $\langle m_\sigma,v_i\rangle=-a_i$。若对所有 $\sigma\in\Sigma$ 和 $1\leq i\leq r$ 均有 $\langle m_\sigma,v_i\rangle\geq-a_i$，则称 $a$ 为\emph{凸的}。

\SourcePage{62}
定义
\begin{equation}\label{eq:cpl-cone}
\operatorname{cpl}(\Sigma)=
\{a\in A^+_{n-1}(X)\otimes\mathbb R:a\text{ 是凸的}\}.
\end{equation}
由于 $X$ 射影且单纯，$\operatorname{cpl}(\Sigma)$ 是一个 $(r-n)$ 维凸锥。而且，$a$ 属于 $\operatorname{cpl}(\Sigma)$ 的内部，当且仅当 $a$ 严格凸，即只要 $\sigma$ 为 $n$ 维且 $v_i\notin\sigma$，就有 $\langle m_\sigma,v_i\rangle>-a_i$，参见 \SourceCite{OdP}。于是可作如下描述。

\begin{proposition}\label{prop:toric-kahler-cone}
若 $X=X_\Sigma$ 射影且单纯，则 $X$ 的 Kähler 锥等于式~\eqref{eq:cpl-cone} 所定义的锥 $\operatorname{cpl}(\Sigma)$ 的内部。
\end{proposition}
\begin{proof}
由 \SourceCite{Baily2}，Kähler 锥包含在 $\operatorname{cpl}(\Sigma)$ 的内部；反向包含关系由 \SourceCite[第 4 节]{AGM1} 证明。
\end{proof}

由于 $\operatorname{cpl}(\Sigma)$ 是强凸多面锥，命题~\ref{prop:toric-kahler-cone} 表明，射影单纯环面簇的 Kähler 锥具有特别良好的结构。

\subsection{Mori 锥}\label{subsec:toric-mori-cone}

若 $X_\Sigma$ 单纯且完备，它的\emph{Mori 锥} $M(X_\Sigma)$ 是
$A_1(X_\Sigma)\otimes\mathbb R\simeq H_2(X_\Sigma,\mathbb R)$
中有效一维圈组成的锥。有时也记为
$NE(X_\Sigma)_{\mathbb R}$ 或 $\overline{NE}(X_\Sigma)_{\mathbb R}$。
在射影情形，我们关注 $M(X_\Sigma)$，是因为它与 Kähler 锥对偶。因此，由命题~\ref{prop:toric-kahler-cone}，$\operatorname{cpl}(\Sigma)$ 与 $M(X_\Sigma)$ 互为对偶。

正如所料，可以用扇 $\Sigma$ 的环面数据显式描述 $M(X_\Sigma)$。为此，取 $(n-1)$ 维锥 $\sigma\in\Sigma(n-1)$；相应的轨道闭包 $C_\sigma\subset X_\Sigma$ 是一条曲线。根据 Reid 的一个观察 \SourceCite[命题 1.6]{Reid3}，$M(X_\Sigma)$ 由这些 $C_\sigma$ 生成。

可以把这一描述写得更明确。令 $v_1,\ldots,v_{n-1}$ 为 $\sigma$ 的整生成元。注意 $\sigma$ 恰好包含在两个 $n$ 维锥中：第一个由 $\sigma$ 与另一个本原整生成元 $v_n$ 张成；第二个则由 $\sigma$ 与另一生成元 $v_{n+1}$ 张成。因此有线性关系
\begin{equation}\label{eq:mori-wall-relation}
\sum_{i=1}^{n+1}\lambda_i v_i=0.
\end{equation}
由于 $v_n$、$v_{n+1}$ 分居 $\sigma$ 两侧，可设 $\lambda_1,\ldots,\lambda_{n+1}$ 为互素整数，且 $\lambda_n,\lambda_{n+1}>0$。于是 $\lambda_i$ 在下列空间中给出一个关系 $\lambda_\sigma$：
\begin{equation}\label{eq:toric-relation-space}
\Lambda_{\mathbb Q}=
\left\{(\lambda_\rho)_{\rho\in\Sigma(1)}:
\lambda_\rho\in\mathbb Q,\ 
\sum_\rho\lambda_\rho v_\rho=0\right\}.
\end{equation}
对式~\eqref{eq:toric-chow-sequence} 取对偶并张量 $\mathbb Q$，得到自然同构
\[
A_1(X_\Sigma)\otimes\mathbb Q\simeq\Lambda_{\mathbb Q}.
\]
在这个同构下，曲线 $C_\sigma\in A_1(X_\Sigma)\otimes\mathbb Q$ 与关系 $\lambda_\sigma\in\Lambda_{\mathbb Q}$ 的联系如下。

\begin{lemma}\label{lem:mori-wall-curve}
设 $\sigma\in\Sigma(n-1)$。那么存在正常数 $c_\sigma$，使上述同构将 $C_\sigma$ 对应到 $c_\sigma\lambda_\sigma$；这里 $\lambda_\sigma$ 是由式~\eqref{eq:mori-wall-relation} 确定的关系。
\end{lemma}

\SourcePage{63}
\begin{proof}
首先，交积使 $A_1(X_\Sigma)\otimes\mathbb Q$ 与 $A_{n-1}(X_\Sigma)\otimes\mathbb Q$ 对偶；而借助 $\mathbb Q^{\Sigma(1)}$ 上的点积以及式~\eqref{eq:toric-chow-sequence}、\eqref{eq:toric-relation-space}，$\Lambda_{\mathbb Q}$ 也与 $A_{n-1}(X_\Sigma)\otimes\mathbb Q$ 对偶。由于同构
$A_1(X_\Sigma)\otimes\mathbb Q\simeq\Lambda_{\mathbb Q}$
与这两个对偶相容，只须找到 $c_\sigma>0$，使对所有 $\rho\in\Sigma(1)$ 有
\[
D_\rho\cdot C_\sigma=c_\sigma D_\rho\cdot\lambda_\sigma.
\]
若 $\rho$ 不属于 $v_i$ 中的任何一个，易见等式两边都是 $0$。若 $\rho$ 是某个 $v_i$，则如通常那样，将 $\rho$ 与 $\rho\cap N$ 的唯一生成元等同，此时 $D_{v_i}\cdot\lambda_\sigma=\lambda_i$。而 \SourceCite[命题 2.7]{Reid3} 表明
\begin{equation}\label{eq:reid-wall-intersection}
\lambda_j D_{v_i}\cdot C_\sigma
=\lambda_i D_{v_j}\cdot C_\sigma,
\qquad 1\leq i,j\leq n+1.
\end{equation}
已知 $\lambda_{n+1}>0$；又因 $\sigma$ 与 $v_{n+1}$ 确定 $\Sigma$ 的一个 $n$ 维锥，$D_{v_{n+1}}\cap C_\sigma$ 只含一个点，即该 $n$ 维锥所对应的不动点。因此
$D_{v_{n+1}}\cdot C_\sigma>0$，从而
$c_\sigma=(D_{v_{n+1}}\cdot C_\sigma)/\lambda_{n+1}>0$。
用式~\eqref{eq:reid-wall-intersection} 很容易验证 $c_\sigma$ 具有所需性质。
\end{proof}

注意，若 $X_\Sigma$ 在 $C_\sigma$ 的某个邻域光滑，则引理~\ref{lem:mori-wall-curve} 中的 $c_\sigma=1$，见 \SourceCite[第 99 页]{Fulton3}。

引理~\ref{lem:mori-wall-curve} 的一个有用推论，是下面的丰沛性判别准则。
\begin{corollary}\label{cor:toric-ample-wall}
除子 $D=\sum_{\rho\in\Sigma(1)}a_\rho D_\rho$ 丰沛，当且仅当它是 Cartier 除子，且对每个 $\sigma\in\Sigma(n-1)$，均满足
$(a_\rho)_{\rho\in\Sigma(1)}\cdot\lambda_\sigma>0$。
\end{corollary}
\begin{proof}
设 $D$ 是 Cartier 除子。由于 $C_\sigma$ 生成 Mori 锥，$D$ 丰沛当且仅当对所有 $\sigma\in\Sigma(n-1)$ 有 $D\cdot C_\sigma>0$。由引理~\ref{lem:mori-wall-curve} 立即得到结论。
\end{proof}

若 $X_\Sigma$ 光滑，Mori 锥 $M(X_\Sigma)$ 还有另一种描述，它涉及\hyperref[subsec:toric-from-homogeneous]{第 3.2.4 小节}定义的\IndexLocation{indexloc:1088:63}本原集合。详情见 \SourceCite{Batyrev1}，相应的丰沛性判别准则见 \SourceCite{Batyrev2}。\index{Mori锥@Mori 锥|hyperpage}

\subsection{辛几何}\label{subsec:toric-symplectic}

辛几何在镜像对称中起重要作用。这里回顾一些基本定义；\SourceCite{Audin} 是一本很好的辛几何参考书。

\emph{辛流形}是赋有闭的、非退化二形式 $\omega$ 的实流形。辛结构把函数转化为向量场：若 $f$ 是流形上的 $C^\infty$ 函数，则流形上存在唯一的向量场 $X_f$，使任意向量场 $X$ 都满足 $\omega(X,X_f)=X(f)$。我们称 $X_f$ 为 $f$ 的\emph{Hamilton 向量场}。\index{辛流形|hyperpage}

对我们而言，最重要的例子是赋有辛形式
$\omega=\sum_{i=1}^{r}dx_i\wedge dy_i$
的 $\mathbb C^r$，其中 $z_i=x_i+iy_i$。$\mathbb C^r$ 的一个特殊性质是，$U(1)^r=(S^1)^r$ 在其上的自然作用是辛的，即对所有 $g\in U(1)^r$ 都有 $g^*\omega=\omega$。

事实上，$U(1)^r$ 的作用还是 Hamilton 作用，含义如下。$U(1)^r$ 的 Lie 代数是 $\mathbb R^r$。元素 $\lambda=(\lambda_1,\ldots,\lambda_r)\in\mathbb R^r$ 在指数映射下对应
$\exp(i\lambda)=(\exp(i\lambda_1),\ldots,\exp(i\lambda_r))\in U(1)^r$。
给定 $\lambda$，在 $\mathbb C^r$ 上得到流
$v\mapsto\exp(it\lambda)\cdot v$，其向量场为
\[
X_\lambda=\sum_{i=1}^{r}\lambda_i
\left(-y_i\frac{\partial}{\partial x_i}
+x_i\frac{\partial}{\partial y_i}\right).
\]

\SourcePage{64}
所谓 Hamilton 作用，就是说对 Lie 代数中的每个 $\lambda$，向量场 $X_\lambda$ 都是 $\mathbb C^r$ 上某个函数的 Hamilton 向量场。

我们用\emph{矩映射}
\[
\mu:\mathbb C^r\longrightarrow(\mathbb R^r)^*
\]
来证明这一点。它定义为
$\mu(z_1,\ldots,z_r)=\tfrac12(|z_1|^2,\ldots,|z_r|^2)$；这里使用的是与 $\mathbb R^r$ 的标准基对偶的 $(\mathbb R^r)^*$ 的基。关键在于，$\lambda=(\lambda_1,\ldots,\lambda_r)\in\mathbb R^r$ 可视为映射 $\lambda:(\mathbb R^r)^*\to\mathbb R$，将它与矩映射复合就得到 $\lambda\circ\mu:\mathbb C^r\to\mathbb R$。简单计算可得
\[
X_\lambda=X_{\lambda\circ\mu},
\]
这便证明了上述作用是 Hamilton 作用。\index{矩映射|hyperpage}

一般而言，矩映射与 Hamilton 作用在辛几何中起重要作用。\hyperref[chap:gromov-witten]{第~\ref*{chap:gromov-witten} 章}给出 Gromov--Witten 不变量的辛几何定义时，我们还会讨论辛几何的其他方面。

\subsection{通过辛约化构造环面簇}\label{subsec:toric-reduction}

下面说明环面簇如何通过称为\emph{辛约化}的过程与辛几何联系起来。设 $X=X_\Sigma$ 由 $N_{\mathbb R}\simeq\mathbb R^n$ 中的扇 $\Sigma$ 确定。如\hyperref[subsec:toric-from-homogeneous]{第 3.2.3 小节}所述，有群 $G=\operatorname{Hom}_{\mathbb Z}(A_{n-1}(X),\mathbb C^*)$。$G$ 的极大紧子群是
\[
G_{\mathbb R}=\operatorname{Hom}_{\mathbb Z}(A_{n-1}(X),U(1)),
\]
它的 Lie 代数为 $\mathfrak g_{\mathbb R}=\operatorname{Hom}_{\mathbb Z}(A_{n-1}(X),\mathbb R)$。因此 $\mathfrak g_{\mathbb R}^*=A_{n-1}(X)\otimes\mathbb R$。

令 $r=|\Sigma(1)|$，则包含映射 $G\subset(\mathbb C^*)^r$ 诱导 $G_{\mathbb R}\subset U(1)^r$。由此得到 $G_{\mathbb R}$ 在 $\mathbb C^r$ 上的 Hamilton 作用，其矩映射显然是复合映射
\[
\mu_\Sigma:\mathbb C^r\xrightarrow{\mu}(\mathbb R^r)^*
\xrightarrow{p}\mathfrak g_{\mathbb R}^*
=A_{n-1}(X)\otimes\mathbb R\simeq\mathbb R^{r-n}.
\]
这里 $\mu$ 是 $U(1)^r$ 的矩映射，$p$ 来自正合列
\begin{equation}\label{eq:moment-exact-sequence}
0\longrightarrow M_{\mathbb R}\longrightarrow(\mathbb R^r)^*
\xrightarrow{p}A_{n-1}(X)\otimes\mathbb R\longrightarrow0,
\end{equation}
该正合列由式~\eqref{eq:toric-chow-sequence} 与 $\mathbb R$ 作张量积得到。映射 $\mu_\Sigma$ 很容易显式描述。首先注意，式~\eqref{eq:moment-exact-sequence} 中的 $(b_1,\ldots,b_r)\in\mathbb R^r$ 在 $M_{\mathbb R}$ 的像上为零，当且仅当在 $N_{\mathbb R}$ 中有 $\sum_{i=1}^r b_iv_i=0$。所有这样的 $(b_1,\ldots,b_r)$ 构成维数为 $r-n$ 的子空间 $\Lambda_{\mathbb R}\subset\mathbb R^r$；注意 $\Lambda_{\mathbb R}=\Lambda_{\mathbb Q}\otimes\mathbb R$，其中 $\Lambda_{\mathbb Q}$ 见式~\eqref{eq:toric-relation-space}。选取 $\Lambda_{\mathbb R}$ 的一组基 $(b_{1,j},\ldots,b_{r,j})$，$1\leq j\leq r-n$，就得到同构 $A_{n-1}(X)\otimes\mathbb R\simeq\mathbb R^{r-n}$，在此同构下 $p$ 变为
\[
p(x_1,\ldots,x_r)=
\left(\sum_{i=1}^r b_{i,1}x_i,\ldots,\sum_{i=1}^r b_{i,r-n}x_i\right).
\]
因此 $\mu_\Sigma$ 的公式为
\begin{equation}\label{eq:toric-moment-map}
\mu_\Sigma(z_1,\ldots,z_r)=\frac12
\left(\sum_{i=1}^r b_{i,1}|z_i|^2,\ldots,
\sum_{i=1}^r b_{i,r-n}|z_i|^2\right).
\end{equation}
我们也看到，$\mu_\Sigma$ 在 $G_{\mathbb R}$ 的每个轨道上取常值。

当 $X$ 射影且单纯时，前面已看到 $A_{n-1}(X)\otimes\mathbb R\simeq H^{1,1}(X,\mathbb R)$。由命题~\ref{prop:toric-kahler-cone}，$X$ 的 Kähler 锥可与 $\operatorname{cpl}(\Sigma)\subset A_{n-1}(X)\otimes\mathbb R$ 的内部等同。于是可如下构造 $X$。

\begin{theorem}\label{thm:toric-symplectic-reduction}
若环面簇 $X=X_\Sigma$ 射影且单纯，而 $a\in A_{n-1}(X)\otimes\mathbb R$ 是 Kähler 类，则 $\mu_\Sigma^{-1}(a)\subset\mathbb C^r-Z(\Sigma)$，且自然映射
\[
\mu_\Sigma^{-1}(a)/G_{\mathbb R}
\longrightarrow(\mathbb C^r-Z(\Sigma))/G=X
\]
\SourcePage{65}
是轨形微分同胚。此外，将 $\mathbb C^r$ 上的辛形式 $\omega$ 限制到 $\mu_\Sigma^{-1}(a)$ 后，它下降为 $\mu_\Sigma^{-1}(a)/G_{\mathbb R}$ 上的辛形式；其上同调类通过上述微分同胚对应于 $a\in H^2(X,\mathbb R)$。
\end{theorem}
\begin{proof}
在定理的陈述中，要注意 $\omega$ 限制到 $\mu_\Sigma^{-1}(a)$ 后并不是辛形式。取商恰好使它变为非退化的。这个过程称为辛约化。由于 $G_{\mathbb R}$ 可以有有限稳定子群，商 $\mu_\Sigma^{-1}(a)/G_{\mathbb R}$ 自然带有辛轨形结构，见\hyperref[app:singular-varieties]{附录 A}。

采用式~\eqref{eq:cpl-cone} 的记号，将 $a\in A_{n-1}(X)\otimes\mathbb R$ 写为 $a=\sum_{i=1}^r a_i[D_i]$。因为 $a$ 是 Kähler 类，命题~\ref{prop:toric-kahler-cone} 表明 $a$ 位于 $\operatorname{cpl}(\Sigma)$ 的内部，即 $a$ 严格凸。考虑由不等式 $\langle m,v_i\rangle\geq-a_i$ 定义的凸集 $\Delta\subset M_{\mathbb R}$。虽然 $\Delta$ 未必是整的，甚至未必是有理的，关键在于 $a$ 的严格凸性意味着多胞形 $\Delta$ 与 $\Sigma$ 在\emph{组合意义下对偶}。也就是说，当 $\sigma$ 遍历 $\Sigma$ 中的锥时，$\Delta$ 的各真面由
\[
\{m\in\Delta:\langle m,v_i\rangle=-a_i
\text{ 对所有 }v_i\in\sigma\}
\]
给出。当 $a_i$ 为整数时，这是\hyperref[subsec:toric-polytopes]{第 3.2.1 小节}讨论的丰沛性的一项性质。$a_i$ 为有理数的情形容易推出，再用连续性论证即可证明一般情形。

知道 $X$ 的扇在组合意义下与 $\Delta$ 对偶之后，就可用 \SourceCite[附录 1，定理 1.4]{Guillemin} 得出
\[
\mu_\Sigma^{-1}(a)/G_{\mathbb R}\simeq X.
\]
此外，定理的最后一个断言来自 \SourceCite[附录 2，式 (1.6)]{Guillemin}；Guillemin 使用的 $\lambda_i$ 在我们的记号下为 $-a_i$。Guillemin 的论证是在光滑情形下给出的，不过稍作修改也适用于 $X$ 单纯的情形。
\end{proof}

利用辛约化构造环面簇，在物理中有一些有趣的推论。特别地，某些称为\emph{规范线性 $\sigma$ 模型}的物理理论，以通过辛约化描述的环面簇为出发点。详情参见\hyperref[sec:gauged-linear-sigma-models]{附录 B.5}。\index{辛约化|hyperpage}\index{规范线性sigma模型@规范线性 $\sigma$ 模型|hyperpage}

\section{GKZ 分解}\label{sec:gkz-decomposition}

\hyperref[chap:moduli]{第~\ref*{chap:moduli} 章}研究 Kähler 模时，我们会看到需要扩大\hyperref[chap:introduction]{第 1 章}定义的 Kähler 模空间。大致的想法是：由一次翻越联系起来的两个 Calabi--Yau 流形，其 Kähler 锥是“相邻”的；这使我们能够将相应的 Kähler 模空间粘合起来。

在环面几何中，“相邻”的 Kähler 锥这一想法自然出现在 \emph{GKZ 分解}中。为了解它的由来，注意由式~\eqref{eq:toric-chow-sequence}，环面簇 $X=X_\Sigma$ 的 Chow 群 $A_{n-1}(X)$ 只依赖于 $\Sigma$ 的一维锥集合 $\Sigma(1)$。因此，具有相同一维锥的两个环面簇，其除子的 Chow 群典范同构。于是可以问：它们各自的 Kähler 锥在与 $\mathbb R$ 作张量积后的 Chow 群中处于什么位置？\index{GKZ分解@GKZ 分解|hyperpage}

\subsection{分解}\label{subsec:gkz-cones}

为精确表述这一想法，在 $N_{\mathbb R}$ 中固定由强凸有理一维锥构成的有限集合 $\Xi$，令 $\Sigma$ 是满足 $\Sigma(1)=\Xi$ 的扇。因为 $A_{n-1}(X)\otimes\mathbb R$ 只依赖于 $\Xi$，将此向量空间记为 $A(\Xi)$。类似地，有效除子类构成的锥 $A^+_{n-1}(X)\otimes\mathbb R$ 也只依赖于 $\Xi$，将它记为 $A^+(\Xi)$。因此 $A^+(\Xi)$ 是 $A(\Xi)\simeq\mathbb R^{r-n}$ 中的锥，其中 $r=|\Xi|$。

\SourcePage{66}
若 $\Sigma$ 是满足 $\Sigma(1)=\Xi$ 的射影单纯扇，则由命题~\ref{prop:toric-kahler-cone}，其 Kähler 锥的闭包是 $\operatorname{cpl}(\Sigma)$。用上述记号，就得到 $(r-n)$ 维锥 $\operatorname{cpl}(\Sigma)\subset A^+(\Xi)$。

现在让 $\Sigma$ 遍历所有这样的扇，便得到位于 $A^+(\Xi)$ 中的一族锥。直观上或许希望这些锥填满整个 $A^+(\Xi)$，但实际并非如此。为得到余下的部分，必须考虑满足 $\Sigma(1)\subset\Xi$ 的射影单纯扇。此时 $A_{n-1}(X_\Sigma)\otimes\mathbb R$ 可能不同于 $A(\Xi)$，但仍可如下定义锥 $\operatorname{cpl}(\Sigma)\subset A^+(\Xi)$。给定 $a=\sum_{i=1}^r a_i[D_i]\in A^+(\Xi)$，对 $\sigma\in\Sigma$，可找到 $m_\sigma\in M_{\mathbb R}$，使得只要 $v_i\in\sigma$ 且 $v_i$ 生成 $\Sigma(1)$ 的一个元素，就有 $\langle m_\sigma,v_i\rangle=-a_i$。若对所有 $\sigma\in\Sigma$ 及 $1\leq i\leq r$，都有 $\langle m_\sigma,v_i\rangle\geq-a_i$，则称 $a$ 为\emph{$\Sigma$ 凸的}。仿照式~\eqref{eq:cpl-cone}，定义
\[
\operatorname{cpl}(\Sigma)=
\{a\in A^+(\Xi):a\text{ 是 }\Sigma\text{ 凸的}\}.
\]
由 \SourceCite{OdP}，此锥仍有维数 $r-n$。值得注意的是，这些锥以非常规整的方式填满 $A^+(\Xi)$，形成 \SourceCite{OdP} 所称的 \emph{Gelfand--Kapranov--Zelevinsky 分解}，简称 \emph{GKZ 分解}。精确描述如下。

\begin{theorem}\label{thm:gkz-fan}
设 $\Xi$ 是 $N_{\mathbb R}$ 中由强凸有理一维锥构成的有限集合。当 $\Sigma$ 遍历所有满足 $\Sigma(1)\subset\Xi$ 的射影单纯扇时，各锥 $\operatorname{cpl}(\Sigma)$ 及其各面在 $A(\Xi)$ 中组成一个扇，其支集为 $A^+(\Xi)$。
\end{theorem}

\hyperref[sec:toric-examples]{第 3.7 节}将给出一个 GKZ 分解的例子，其中有两个锥 $\operatorname{cpl}(\Sigma)$ 和 $\operatorname{cpl}(\Sigma')$，满足 $\Sigma(1)=\Xi$ 而 $\Sigma'(1)\subsetneq\Xi$。

一般而言，可以问 GKZ 分解中的两个锥何时有一个共同的余维数为一的面。答案很漂亮：当 $N_{\mathbb R}$ 中对应的两个扇通过细分，即增加或减去一条一维锥，或通过一次“翻越”联系起来时，就会出现这种情形；详情见 \SourceCite{OdP}。\hyperref[subsec:enlarged-kahler-moduli]{第~\ref*{subsec:enlarged-kahler-moduli} 小节}将讨论翻越。还应提到，非单纯扇对应于 GKZ 分解中某些维数小于 $r-n$ 的锥。

\subsection{次级扇与 Gale 变换}\label{subsec:secondary-fan}

GKZ 分解可以扩充为一个完备扇，其支集是整个 $A(\Xi)$。具体做法如下。首先，把 $\Xi$ 中每条一维锥与它的本原整生成元相对应，将 $\Xi$ 视为 $N_{\mathbb R}$ 的一个子集。然后定义子集 $\Xi^+\subset N_{\mathbb R}\times\mathbb R$：
\begin{equation}\label{eq:extended-rays}
\Xi^+=(\Xi\cup\{0\})\times\{1\}.
\end{equation}
我们关心的是 $\Xi^+$ 的\emph{次级扇}。其定义类似于 GKZ 分解的定义：将 GKZ 分解中出现的射影单纯扇替换为凸包 $\operatorname{Conv}(\Xi^+)$ 的\emph{\IndexLocation{indexloc:1226:66}正则三角剖分} $\mathcal T$。$\mathcal T$ 中各单形的顶点必须属于 $\Xi^+$；所谓正则，是指 $\mathcal T$ 可由一个\IndexLocation{indexloc:0182:66}严格凸函数确定，正如射影单纯扇 $\Sigma$ 可由严格凸支撑函数 $a\in\operatorname{cpl}(\Sigma)$ 确定一样。与 $\mathcal T$ 对应的是锥 $C(\mathcal T)\subset A(\Xi)$，它由确定同一三角剖分 $\mathcal T$ 的严格凸支撑函数及其极限构成。各 $C(\mathcal T)$ 是次级扇的极大锥；次级扇是 $A(\Xi)$ 中的完备扇。\index{次级扇|hyperpage}\index{Gale变换@Gale 变换|hyperpage}

给定三角剖分 $\mathcal T$，锥 $C(\mathcal T)$ 还有如下描述。由式~\eqref{eq:toric-chow-sequence} 可见，$\Xi^+\subset N_{\mathbb R}\times\mathbb R$ 给出
\begin{equation}\label{eq:gale-exact-sequence}
0\longrightarrow M_{\mathbb R}\oplus\mathbb R
\longrightarrow\mathbb R^{\Xi^+}\longrightarrow A(\Xi)
\longrightarrow0.
\end{equation}

\SourcePage{67}
对每个 $v\in\Xi^+$，令 $e_v^*\in A(\Xi)$ 为 $\mathbb R^{\Xi^+}$ 中对应于 $v$ 的标准基向量在式~\eqref{eq:gale-exact-sequence} 的满射下的像。元素集合 $\{e_v^*:v\in\Xi^+\}$ 称为 $\Xi^+$ 的 \emph{Gale 变换}。读者可以验证，若 $\Sigma(1)=\Xi$ 且 $v=(\rho,1)\in\Xi^+-\{(0,1)\}$，则 $e_v^*=[D_\rho]$，即 $\rho\in\Sigma(1)$ 对应的环面除子在 $A(\Xi)=A_{n-1}(X_\Sigma)\otimes\mathbb R$ 中的类。

设 $\mathcal T$ 是上述 $\operatorname{Conv}(\Xi^+)$ 的正则三角剖分。每个单形 $\sigma\in\mathcal T$ 给出锥 $C(\sigma)\subset A(\Xi)$，其生成元是那些 $v$ 不是 $\sigma$ 顶点的 $e_v^*$。可以证明
\[
C(\mathcal T)=\bigcap_{\sigma\in\mathcal T}C(\sigma).
\]
若三角剖分 $\mathcal T$ 的每个极大单形都以 $(0,1)$ 为顶点，则因为 $(0,1)\in\operatorname{Conv}(\Xi^+)\subset N_{\mathbb R}\times\{1\}\simeq N_{\mathbb R}$ 是内点，从顶点 $(0,1)$ 向 $\mathcal T$ 中各单形的面作锥，便得到一个完备扇 $\Sigma$。所得扇满足 $\Sigma(1)\subset\Xi$，且 $C(\mathcal T)=\operatorname{cpl}(\Sigma)$。这个过程可以反过来进行，因此次级扇含有 GKZ 扇作为子扇。次级扇的更完整描述见 \SourceCite{GKZ2}，它与 GKZ 分解的关系见 \SourceCite{OdP}。

由式~\eqref{eq:gale-exact-sequence} 可见，对偶空间 $A(\Xi)^*$ 为
\[
A(\Xi)^*=
\left\{\lambda=(\lambda_v)\in\mathbb R^{\Xi^+}:
\sum_{v\in\Xi^+}\lambda_v\,v=0\right\}.
\]
这是 $\Xi^+$ 中各元素的线性关系构成的向量空间，类似于式~\eqref{eq:toric-relation-space} 定义的 $\Lambda_{\mathbb Q}$。注意，$A(\Xi)^*$ 中有一个显然的整数格，即系数全为整数的关系构成的格。此外，Gale 变换中的每个 $e_v^*\in A(\Xi)$ 都是 $A(\Xi)^*$ 上的线性泛函，容易算得 $e_v^*(\lambda)=\lambda_v$。\hyperref[chap:moduli]{第~\ref*{chap:moduli} 章}将按以下方式使用这一事实：固定 $A(\Xi)^*$ 的一组基，将每个基向量视为一行 $\lambda=(\lambda_v)$，这样便得到一个矩阵；若在 $A(\Xi)$ 中使用对偶基，则这个矩阵的各列就是各 $e_v^*$。这种描述使 Gale 变换易于处理。

一个特别好的情形是 $|\Xi^+|=n+3$，其中 $n$ 是格 $N$ 和 $M$ 的秩。此时 $A(\Xi)$ 的维数为 $2$，Gale 变换 $\{e_v^*\}$ 是平面上的一组向量。这时容易看到，次级扇就是以 Gale 变换中的各向量所张成的一维锥为一维锥的完备二维扇；\SourceCite{DHSS} 明确指出了这一点。

\subsection{与辛约化的关系}\label{subsec:gkz-reduction}

本节最后考察矩映射与 GKZ 分解的关系。首先注意，矩映射 $\mu_\Sigma$ 只依赖于 $\Sigma(1)$，因为正合列~\eqref{eq:moment-exact-sequence} 由 $\Sigma(1)$ 确定。因此可将矩映射写为 $\mu_\Xi:\mathbb C^\Xi\to A(\Xi)$。它的像是 $A^+(\Xi)$，所以有
\[
\mu_\Xi:\mathbb C^\Xi\longrightarrow A^+(\Xi).
\]
由于作用在 $\mathbb C^\Xi$ 上的群 $G$ 也只依赖于 $\Xi$，将它记为 $G(\Xi)$。

设 GKZ 分解中有一个 $(r-n)$ 维锥 $\operatorname{cpl}(\Sigma)\subset A^+(\Xi)$。取其内部的 $a$，形成商
\[
\mu_\Xi^{-1}(a)/G(\Xi)_{\mathbb R}.
\]
若 $\Sigma(1)=\Xi$ 且 $\Sigma$ 射影、单纯，由定理~\ref{thm:toric-symplectic-reduction}，这个商就是 $X_\Sigma$。当 $\Sigma(1)$ 是 $\Xi$ 的真子集时，上述定理不再适用，但以下推广告诉我们，得到的仍是 $X_\Sigma$。

\SourcePage{68}
\begin{theorem}\label{thm:gkz-symplectic-reduction}
若 $\Sigma$ 是满足 $\Sigma(1)\subset\Xi$ 的射影单纯扇，且 $a\in A^+(\Xi)$ 位于 $\operatorname{cpl}(\Sigma)$ 的内部，则有自然的轨形微分同胚
\[
\mu_\Xi^{-1}(a)/G(\Xi)_{\mathbb R}\simeq X_\Sigma.
\]
\end{theorem}
\begin{proof}
除了群 $G(\Xi)$ 与矩映射 $\mu_\Xi$，还有由 $\Sigma$ 得到的 $G(\Sigma)$ 和 $\mu_\Sigma$。为联系二者，令 $C=\Xi-\Sigma(1)$，于是 $\mathbb Z^\Xi=\mathbb Z^{\Sigma(1)}\oplus\mathbb Z^C$。投影 $\mathbb Z^\Xi\to\mathbb Z^{\Sigma(1)}$ 诱导正合列
\begin{equation}\label{eq:gkz-group-sequence}
1\longrightarrow G(\Sigma)\longrightarrow G(\Xi)
\longrightarrow(\mathbb C^*)^C\longrightarrow1,
\end{equation}
使得相对于分解 $\mathbb C^\Xi=\mathbb C^{\Sigma(1)}\oplus\mathbb C^C$，$G(\Sigma)$ 在第二个因子上的作用是平凡的。此外，可通过以下交换图联系矩映射：
\[
\begin{array}{rccccc}
\mu_\Xi:&\mathbb C^\Xi&\longrightarrow&\mathbb R^\Xi&\longrightarrow&A(\Xi)\\
&\big\downarrow&&\big\downarrow&&\big\downarrow\\
\mu_\Sigma:&\mathbb C^{\Sigma(1)}&\longrightarrow&\mathbb R^{\Sigma(1)}
&\longrightarrow&A_{n-1}(X_\Sigma)\otimes\mathbb R,
\end{array}
\]
其中竖直箭头是自然投影。

若 $a=\sum_{i=1}^r a_i[D_i]$ 位于 $\operatorname{cpl}(\Sigma)$ 的内部，则 \SourceCite[命题 3.1]{OdP} 的证明表明 $a$ 严格凸，即当 $\sigma\in\Sigma$ 为 $n$ 维、$v_i$ 不是 $\sigma$ 的生成元时，有 $\langle m_\sigma,v_i\rangle>-a_i$。令 $a_0$ 为 $a$ 在 $A_{n-1}(X_\Sigma)\otimes\mathbb R$ 中的像，则 $a_0$ 也严格凸。由定理~\ref{thm:toric-symplectic-reduction}，只须证明
\[
\mu_\Sigma^{-1}(a_0)/G(\Sigma)_{\mathbb R}
\simeq\mu_\Xi^{-1}(a)/G(\Xi)_{\mathbb R}.
\]

由于 $\Sigma$ 完备，每个 $v_i\in C$ 都位于某个 $n$ 维锥 $\sigma\in\Sigma$ 中；这里仍按惯例将一维锥与其生成元等同。于是 $v_i=\sum_j c_{ij}v_j$，其中 $c_{ij}\geq0$，求和遍历那些 $v_j$ 为 $\sigma$ 生成元的 $j$；本证明余下部分沿用这一约定。严格凸条件 $\langle m_\sigma,v_i\rangle>-a_i$ 给出不等式
\[
a_i>\sum_j c_{ij}a_j,\qquad v_i\in C.
\]
设 $\mathbb C^\Xi$ 的坐标为 $z_1,\ldots,z_r$，选择编号使 $v_1,\ldots,v_k\in\Sigma(1)$，$v_{k+1},\ldots,v_r\in C$。不难证明，$(z_1,\ldots,z_r)\in\mu_\Xi^{-1}(a)$，当且仅当 $(z_1,\ldots,z_k)\in\mu_\Sigma^{-1}(a_0)$，并且
\begin{equation}\label{eq:gkz-extra-coordinates}
\frac12|z_i|^2=\frac12\sum_j c_{ij}|z_j|^2
+a_i-\sum_j c_{ij}a_j,
\qquad k+1\leq i\leq r.
\end{equation}

定义映射 $\phi:\mu_\Sigma^{-1}(a_0)\to\mu_\Xi^{-1}(a)$：将 $(z_1,\ldots,z_k)\in\mu_\Sigma^{-1}(a_0)$ 送到 $(z_1,\ldots,z_r)\in\mu_\Xi^{-1}(a)$，其中 $z_{k+1},\ldots,z_r$ 取为满足式~\eqref{eq:gkz-extra-coordinates} 的正实数；根据上述不等式，该式右边为正。于是 $\phi$ 对 $G(\Sigma)_{\mathbb R}$ 与 $G(\Xi)_{\mathbb R}$ 的作用是等变的，而且由式~\eqref{eq:gkz-group-sequence}，$G(\Xi)_{\mathbb R}$ 的每个轨道都与 $\phi$ 的像相交。最后，设 $z,w\in\mu_\Sigma^{-1}(a_0)$、$g\in G(\Xi)_{\mathbb R}$ 满足 $g\phi(z)=\phi(w)$。由于 $g$ 不改变绝对值，式~\eqref{eq:gkz-extra-coordinates} 表明 $g$ 在 $z_{k+1},\ldots,z_r$ 上的作用平凡。这些数又都非零，因此由式~\eqref{eq:gkz-group-sequence} 得 $g\in G(\Sigma)_{\mathbb R}$。所以 $\phi$ 在轨道之间诱导双射，容易看出它是轨形微分同胚。
\end{proof}

定理~\ref{thm:gkz-symplectic-reduction} 表明，矩映射 $\mu_\Xi$ 不仅能构造一个环面簇，而且能构造所有满足 $\Sigma(1)\subset\Xi$ 的射影单纯环面簇 $X_\Sigma$。这与\emph{\IndexLocation{indexloc:0924:68}多重镜像}现象有关，第\hyperref[chap:mirror-constructions]{4 章}将对此展开讨论。\SourceCite{Witten5} 的\IndexLocation{indexloc:0354:68}规范线性 $\sigma$ 模型也与此相关。
\SourcePage{69}
在这些模型中，元素 $a\in A(\Xi)$ 是一个参数；$a$ 所处的位置不同，就得到不同的物理理论。若 $a$ 位于某个锥 $\operatorname{cpl}(\Sigma)$ 的内部，则理论涉及相应的环面簇 $X_\Sigma$；若它位于 $A^+(\Xi)$ 之外，则得到很不一样的理论，例如 Landau--Ginzburg 理论。物理文献将这种现象称为\emph{相}，例如参见 \SourceCite{MP1}。

\section{Fano 簇与自反多胞形}\label{sec:fano-reflexive}

射影空间的反典范类是丰沛的。更一般地，任何具有此性质的光滑完备簇都称为 \emph{Fano 簇}。为了研究镜像对称，还需考虑有奇点的 Fano 簇。为给出定义，回顾\hyperref[app:singular-varieties]{附录 A}：Cohen--Macaulay 簇 $V$ 有对偶化层 $\widehat\Omega_V^n$，其中 $n=\dim V$；而 $\widehat\Omega_V^n$ 是线丛，当且仅当 $V$ 是 Gorenstein 的。

Fano 性意味着对偶化层的对偶是丰沛的。特别地，这说明我们可能希望 $\widehat\Omega_V^n$ 是线丛，也就是希望这样的簇为 Gorenstein 簇。因此有如下定义。
\begin{definition}\label{def:fano}
若完备的 $n$ 维 Gorenstein 簇 $V$ 的线丛 $\widehat\Omega_V^n$ 的对偶是丰沛的，则称 $V$ 为 \emph{Fano 簇}。
\end{definition}

对环面簇，可以如下刻画这些概念。首先回顾，任意环面簇 $X$ 上的对偶化层都有简单描述
\[
\widehat\Omega_X^n=\mathcal O_X\left(-\sum_\rho D_\rho\right),
\]
见 \SourceCite{Danilov,Fulton3,Oda}。用典范除子 $K_X$ 表示，就是 $K_X=-\sum_\rho D_\rho$。因此，$X$ 是 Gorenstein 的，当且仅当 $\sum_\rho D_\rho$ 是 Cartier 除子。于是可如下刻画 Fano 环面簇。
\begin{lemma}\label{lem:toric-fano}
完备的 $n$ 维环面簇 $X$ 是 Fano 簇，当且仅当 $\sum_\rho D_\rho$ 是丰沛 Cartier 除子。
\end{lemma}

由于每个丰沛除子都来自一个多胞形，Fano 环面簇上的丰沛除子 $-K_X$ 确定了一个多胞形 $\Delta$，它具有一些很特殊的性质。这就引出了 Batyrev 的\emph{自反多胞形}概念 \SourceCite{Batyrev4}。精确定义如下。
\begin{definition}\label{def:reflexive-polytope}
若 $n$ 维整多胞形 $\Delta\subset M_{\mathbb R}\simeq\mathbb R^n$ 满足以下两个条件，则称 $\Delta$ 为\emph{自反}的：
\begin{enumerate}[label=(\roman*)]
\item $\Delta$ 的每个刻面 $\Gamma$ 都位于形如 $\{m\in M_{\mathbb R}:\langle m,v_\Gamma\rangle=-1\}$ 的支撑仿射超平面上，其中 $v_\Gamma\in N$；
\item $\operatorname{Int}(\Delta)\cap M=\{0\}$。
\end{enumerate}
\end{definition}

自反多胞形具有非常漂亮的组合对偶性。设 $\Delta$ 为整多胞形，$\Delta^\circ$ 是\hyperref[subsec:toric-from-polytopes]{第 3.2.2 小节}定义的极多胞形。除了 $(\Delta^\circ)^\circ=\Delta$ 之外，\SourceCite{Batyrev4} 还证明，$\Delta$ 与 $\Delta^\circ$ 的基本对偶性如下。
\begin{lemma}\label{lem:reflexive-duality}
$\Delta$ 自反，当且仅当 $\Delta^\circ$ 自反。
\end{lemma}

自反多胞形在这里之所以有趣，是因为以下结果刻画了何时 $\mathbb P_\Delta$ 是 Fano 簇。
\begin{proposition}\label{prop:reflexive-fano}
$\Delta$ 自反，当且仅当 $\mathbb P_\Delta$ 是 Fano 簇，且 $\Delta$ 是其反典范除子对应的多胞形。
\end{proposition}

\SourcePage{70}
\begin{proof}
设 $\Delta$ 自反，则定义~\ref{def:reflexive-polytope} 中的 $v_\Gamma$ 恰好是 $\mathbb P_\Delta$ 的扇的一维锥生成元。容易看出，支撑函数 $\phi_\Delta$ 在每个 $v_\Gamma$ 上都取值 $-1$。因此，$\mathbb P_\Delta$ 上相应的丰沛除子是 $\sum_\rho D_\rho=-K_{\mathbb P_\Delta}$，所以 $\mathbb P_\Delta$ 是 Fano 簇。逆命题也同样容易证明。
\end{proof}

Fano 环面簇最简单的例子是 $\mathbb P^n$。接着考虑加权射影空间，此时答案稍微有趣一些。
\begin{lemma}\label{lem:weighted-projective-fano}
设 $X=\mathbb P(q_0,\ldots,q_n)$ 是使用归一化权重的加权射影空间，令 $q=\sum_{i=0}^n q_i$。则 $X$ 是 Fano 簇，当且仅当对每个 $i$ 都有 $q_i\mid q$。
\end{lemma}
\begin{proof}
$X=\mathbb P(q_0,\ldots,q_n)$ 的扇的一维锥生成元为 $v_0,\ldots,v_n$，在 $n$ 维格 $N$ 中满足 $\sum_{i=0}^n q_iv_i=0$。扇的极大锥由 $\{v_0,\ldots,v_n\}$ 的各 $n$ 元子集生成。首先判定除子 $D=\sum_{i=0}^n D_i$ 何时是 Cartier 除子。

对每个 $i$，容易看出存在唯一的 $m_i\in M\otimes\mathbb Q$，使对所有 $j\ne i$ 都有 $\langle m_i,v_j\rangle=-1$。于是 $D$ 是 Cartier 除子，当且仅当每个 $m_i$ 都属于 $M$。不过，关系 $\sum_{i=0}^n q_iv_i=0$ 意味着
\[
\langle m_i,v_i\rangle=\frac{\sum_{j\ne i}q_j}{q_i}
=\frac q{q_i}-1.
\]
由于 $v_0,\ldots,v_n$ 生成 $N$，可见 $m_i\in M$ 当且仅当 $q_i\mid q$；对所有 $i$ 同时要求这一条件，就得到上述判别。一旦知道 $D$ 是 Cartier 除子，\hyperref[sec:cones-fans]{第 3.1 节}的丰沛性准则便化为不等式 $\langle m_i,v_i\rangle>-1$，这在此处显然成立。
\end{proof}

\hyperref[sec:toric-examples]{第 3.7 节}将给出更多 Fano 环面簇的例子。在任意给定维数中，模去幺模变换后，自反多胞形只有有限多个；这意味着 $n$ 维 Fano 环面簇的同构类也只有有限多个，见 \SourceCite{Batyrev4}。低维光滑 Fano 环面簇已得到分类，见 \SourceCite{Batyrev6,Oda}；同时，人们正在尝试分类所有四维自反多胞形，见 \SourceCite{KS3,Skarke}。

由于自反多胞形成对出现为 $\Delta$、$\Delta^\circ$，我们得到在某种意义下对偶的 Fano 环面簇 $\mathbb P_\Delta$、$\mathbb P_{\Delta^\circ}$。\hyperref[chap:mirror-constructions]{第 4 章}将利用这些环面簇构造 Calabi--Yau 超曲面的“对偶”族，它们在镜像对称中十分重要。\index{Fano簇@Fano 簇|hyperpage}

\section{环面簇的自同构}\label{sec:toric-automorphisms}

本节描述完备单纯环面簇 $X$ 的自同构群。这些结果源于 Demazure \SourceCite{Demazure}，他在 $X$ 光滑时计算了 $\operatorname{Aut}(X)$。\SourceCite{Cox} 证明，Demazure 对 $\operatorname{Aut}(X)$ 的描述在单纯情形同样成立。

若 $X$ 由扇 $\Sigma$ 给出，则定理~\ref{thm:toric-quotient} 给出商表示 $X=(\mathbb C^{\Sigma(1)}-Z(\Sigma))/G$，其中 $G=\operatorname{Hom}_{\mathbb Z}(A_{n-1}(X),\mathbb C^*)$。还需回顾齐次坐标环 $S=\mathbb C[x_\rho:\rho\in\Sigma(1)]$，其分次由 $A_{n-1}(X)$ 给出。

我们将通过寻找与 $G$ 的作用交换的 $\mathbb C^{\Sigma(1)}-Z(\Sigma)$ 的自同构，来构造 $X=(\mathbb C^{\Sigma(1)}-Z(\Sigma))/G$ 的自同构。一类显然的自同构由“大环面” $(\mathbb C^*)^{\Sigma(1)}$ 中的元素给出。正合列~\eqref{eq:toric-g-sequence}
\[
1\longrightarrow G\longrightarrow(\mathbb C^*)^{\Sigma(1)}
\longrightarrow T_N\longrightarrow1
\]
\SourcePage{71}
表明，这一类给出来自环面的 $X$ 的自同构，即 $T_N\subset\operatorname{Aut}(X)$。

下一类自同构来自\emph{根}。$X$ 的一个根是一个有序对 $(x_\rho,x^D)$，其中 $x_\rho$ 是我们的一个变量，$x^D\in S$ 是不同于 $x_\rho$、但与它同次的单项式，即在 $A_{n-1}(X)$ 中有 $\deg(x^D)=\deg(x_\rho)$。注意，$x_\rho$ 不可能整除 $x^D$，否则 $x^D/x_\rho$ 就是一个次数为 $0$ 的非平凡单项式；由于 $X$ 完备，这样的单项式不存在。

根 $(x_\rho,x^D)$ 如下确定 $\mathbb C^{\Sigma(1)}$ 的自同构。将 $\mathbb C^{\Sigma(1)}$ 的一点写为 $(x_\rho,\mathbf x)$，其中 $\mathbf x$ 是以 $\Sigma(1)-\{\rho\}$ 为指标的向量。于是得到单参数自同构族
\begin{equation}\label{eq:toric-root-automorphism}
y_\lambda(x_\rho,\mathbf x)=
(x_\rho+\lambda\mathbf x^D,\mathbf x),
\end{equation}
其中 $\lambda\in\mathbb C$，$\mathbf x^D$ 是单项式 $x^D$ 在 $\mathbf x$ 处的值；由于 $x_\rho$ 不整除 $x^D$，这个取值有意义。可以证明，$y_\lambda$ 是 $\mathbb C^{\Sigma(1)}-Z(\Sigma)$ 的自同构，并与 $G$ 的作用交换，因此下降为 $\operatorname{Aut}(X)$ 中的元素。

最后一类自同构来自扇 $\Sigma$ 的对称，即置换 $\Sigma$ 中各锥的 $N$ 的自同构。特别地，它们置换一维锥，因此给出 $\mathbb C^{\Sigma(1)}-Z(\Sigma)$ 的自同构；这些自同构与 $G$ 的作用交换，因而也下降到 $X$。

\begin{theorem}\label{thm:toric-automorphism-group}
若 $X$ 是单纯完备环面簇，则来自环面、根以及扇对称的上述三类自同构生成 $\operatorname{Aut}(X)$。此外，来自环面与根的前两类生成 $\operatorname{Aut}(X)$ 的单位元连通分支，而且
\[
\begin{aligned}
\dim\operatorname{Aut}(X)
&=n+\text{根的个数}\\
&=n+\sum_\rho\bigl(\dim S_{\deg(x_\rho)}-1\bigr).
\end{aligned}
\]
\end{theorem}

例如，$\mathbb P^n$ 的\IndexLocation{indexloc:1369:71}齐次坐标环是通常的分次环 $\mathbb C[x_0,\ldots,x_n]$。这意味着，根恰好是所有 $i\ne j$ 的有序对 $(x_i,x_j)$。这样的有序对有 $(n+1)^2-(n+1)$ 个，因此由定理~\ref{thm:toric-automorphism-group}，$\dim\operatorname{Aut}(\mathbb P^n)=(n+1)^2-1$。这当然就是 $\operatorname{PGL}(n+1,\mathbb C)$ 的维数。文献 \SourceCite{Cox,Demazure,Oda} 给出了 $\operatorname{Aut}(X)$ 作为代数群的结构的更多细节。

下面介绍 \SourceCite{AGM1} 给出的一种计算 $\operatorname{Aut}(X)$ 维数的几何方法。假设\IndexLocation{indexloc:1375:71}单纯环面簇 $X$ 是 Gorenstein 的，则反典范类 $-K_X=\sum_\rho D_\rho$ 是 Cartier 除子，给出多胞形
\[
\Delta=\Delta_{-K_X}=
\{m\in M_{\mathbb R}:\langle m,v_\rho\rangle\geq-1\}.
\]
注意，$0$ 在 $\Delta$ 的内部，因此 $\dim\Delta=n$，不过 $\Delta$ 未必是整的。

\begin{proposition}\label{prop:automorphism-facets}
若 $X$ 单纯、完备且 Gorenstein，而 $\Delta$ 如上定义，则
\[
\dim\operatorname{Aut}(X)=n+\sum_\Gamma l^*(\Gamma),
\]
其中求和遍历 $\Delta$ 的所有刻面 $\Gamma$。
\end{proposition}
\begin{proof}
回顾式~\eqref{eq:lattice-points}，$l^*(\Gamma)$ 是刻面 $\Gamma$ 的相对内部中的格点数。为证明命题，设 $(x_\rho,x^D)$ 是一个根，写 $x^D=\prod_{\rho'\ne\rho}x_{\rho'}^{a_{\rho'}}$。由正合列~\eqref{eq:toric-chow-sequence}，$\deg(x_\rho)=\deg(x^D)$ 当且仅当存在 $\widetilde m\in M$，使
\begin{equation}\label{eq:root-lattice-pairings}
\begin{aligned}
\langle\widetilde m,v_\rho\rangle&=-1,\\
\langle\widetilde m,v_{\rho'}\rangle&=a_{\rho'},\qquad\rho'\ne\rho.
\end{aligned}
\end{equation}

\SourcePage{72}
由于 $\Delta$ 由不等式 $\langle m,v_\rho\rangle\geq-1$ 及 $\rho'\ne\rho$ 时的 $\langle m,v_{\rho'}\rangle\geq-1$ 定义，式~\eqref{eq:root-lattice-pairings} 表明 $\widetilde m\in\Delta$，因为对 $\rho'\ne\rho$ 有 $a_{\rho'}\geq0>-1$。事实上，后面这些严格不等式与 $\langle\widetilde m,v_\rho\rangle=-1$ 意味着，$\widetilde m$ 位于 $\Delta$ 中由 $\langle m,v_\rho\rangle=-1$ 定义的刻面的相对内部。反过来，给定这个刻面相对内部的格点 $\widetilde m$，取
\[
x^D=\prod_{\rho'\ne\rho}x_{\rho'}^{\langle\widetilde m,v_{\rho'}\rangle},
\]
容易验证 $(x_\rho,x^D)$ 是一个根。
\end{proof}

举例说明这种方法：$\mathbb P^4$ 是 Fano 簇，因此 $\mathbb P^4=\mathbb P_\Delta$，其中 $\Delta=\Delta_{-K_{\mathbb P^4}}$ 自反。我们知道自同构群 $\operatorname{Aut}(\mathbb P^4)$，但 $X^\circ=\mathbb P_{\Delta^\circ}$ 的自同构如何？这里 $\Delta^\circ$ 是 $\Delta$ 的对偶多胞形。\hyperref[sec:quintic-revisited]{第 4.2 节}将说明，单形 $\Delta^\circ$ 的顶点为
\[
\begin{gathered}
(-1,-1,-1,-1),\\
(1,0,0,0),\\
(0,1,0,0),\\
(0,0,1,0),\\
(0,0,0,1).
\end{gathered}
\]
注意，这个多胞形的刻面没有内部格点。由于 $\Delta^\circ$ 是 $X^\circ$ 的反典范类对应的多胞形，命题~\ref{prop:automorphism-facets} 表明 $X^\circ$ 没有根。因此，$\operatorname{Aut}(X^\circ)$ 的单位元连通分支恰好是 $X^\circ$ 的环面。特别地，虽然原来的环面簇 $X=\mathbb P^4$ 有很大的自同构群，它的对偶 $X^\circ$ 的自同构群却相对较小。\index{自同构群|hyperpage}

\section{例题}\label{sec:toric-examples}

本节给出三个例题，说明目前已讨论的环面簇理论的不同方面。对不太熟悉环面几何的读者，仔细完成这些例题的细节是一项很好的练习。

\begin{example}\label{ex:toric-projective-space}
环面簇最基本的例子是 $\mathbb P^n$。此时 $N=\mathbb Z^n$，一维锥的生成元是标准基 $e_1,\ldots,e_n$ 以及 $e_0=-\sum_{i=1}^n e_i$；仍按惯例不区分一维锥与其生成元。由 $\{e_0,\ldots,e_n\}$ 的所有真子集所生成的锥组成一个光滑扇。正合列~\eqref{eq:toric-chow-sequence} 变为
\[
0\longrightarrow\mathbb Z^n\longrightarrow\mathbb Z^{n+1}
\longrightarrow\mathbb Z\longrightarrow0,
\]
其中 $(a_0,\ldots,a_n)\in\mathbb Z^{n+1}$ 映到 $a_0+\cdots+a_n$。因此，\IndexLocation{indexloc:1369:72}齐次坐标环就是赋有通常分次的 $\mathbb C[x_0,\ldots,x_n]$，定理~\ref{thm:toric-quotient} 给出的 $\mathbb P^n$ 的构造就是通常的 $\mathbb P^n=(\mathbb C^{n+1}-\{0\})/\mathbb C^*$。

矩映射为 $\mu(z_0,\ldots,z_n)=\tfrac12\sum_{i=0}^n|z_i|^2$，所以对 $r>0$，$\mu^{-1}(r)$ 是一个球面，而定理~\ref{thm:toric-symplectic-reduction} 的构造就是纤维为 $S^1$ 的 Hopf 纤维化 $S^{2n+1}\to\mathbb P^n$。

最后，对应于 $e_0$ 的除子 $D_0$ 是丰沛的，其对应多胞形为单形
\[
\Delta_{D_0}=\{(c_1,\ldots,c_n):c_i\geq0,\ c_1+\cdots+c_n\leq1\}.
\]
\end{example}

\begin{example}\label{ex:weighted-resolution}
下一个例题涉及\IndexLocation{indexloc:1432:72}加权射影空间及其环面消解。在 $N=\mathbb Z^4$ 中考虑向量
\[
v_0=(-1,-2,-2,-2),\quad v_1=e_1,\quad v_2=e_2,
\quad v_3=e_3,\quad v_4=e_4.
\]
\SourcePage{73}
这些向量满足关系 $v_0+v_1+2v_2+2v_3+2v_4=0$。使用由 $\{v_0,v_1,v_2,v_3,v_4\}$ 的真子集构成的扇，得到加权射影空间 $\mathbb P(1,1,2,2,2)$；由引理~\ref{lem:weighted-projective-fano}，它是 Fano 簇。正合列~\eqref{eq:toric-chow-sequence} 为
\[
0\longrightarrow\mathbb Z^4\longrightarrow\mathbb Z^5
\longrightarrow\mathbb Z\longrightarrow0,
\]
其中 $(a_0,a_1,a_2,a_3,a_4)\in\mathbb Z^5$ 映到 $a_0+a_1+2a_2+2a_3+2a_4$。因此，齐次坐标环是 $\mathbb C[x_0,x_1,x_2,x_3,x_4]$，其中 $x_0,x_1$ 的次数为 $1$，$x_2,x_3,x_4$ 的次数为 $2$。

为消解 $\mathbb P(1,1,2,2,2)$ 的奇点，取由 $v_0,v_1$ 生成的奇异锥，加入生成元
\[
v_5=\tfrac12(v_0+v_1)=(0,-1,-1,-1).
\]
这给出新扇 $\Sigma$：原扇中每个包含 $v_0,v_1$ 的四维锥，都被分成 $\Sigma$ 中的两个锥。$\Sigma$ 是光滑的，因此对应的环面簇 $X=X_\Sigma$ 是 $\mathbb P(1,1,2,2,2)$ 的奇点消解。

对 $X$，可以看到群 $G=(\mathbb C^*)^2$ 在 $\mathbb C^{\Sigma(1)}=\mathbb C^6$ 上的作用为
\[
\begin{aligned}
&(\lambda,\mu)(x_0,x_1,x_2,x_3,x_4,x_5)\\
&\qquad=(\lambda x_0,\lambda x_1,\mu x_2,\mu x_3,\mu x_4,
\lambda^{-2}\mu x_5).
\end{aligned}
\]
此外，仅有的两个\IndexLocation{indexloc:1087:73}本原集合是 $\{v_0,v_1\}$ 与 $\{v_2,v_3,v_4,v_5\}$。由式~\eqref{eq:toric-primitive-z}，
$Z(\Sigma)=\{x_0=x_1=0\}\cup\{x_2=x_3=x_4=x_5=0\}$。因此
\[
X=(\mathbb C^6-Z(\Sigma))/(\mathbb C^*)^2.
\]

还可以描述矩映射与 GKZ 分解。因为 $v_0,v_1,v_2,v_3,v_4,v_5$ 之间的关系由 $v_2+v_3+v_4+v_5=0$ 及 $v_0+v_1-2v_5=0$ 生成，式~\eqref{eq:toric-moment-map} 表明矩映射为
\[
\begin{aligned}
&\mu_\Sigma(z_0,z_1,z_2,z_3,z_4,z_5)\\
&\quad=\tfrac12\bigl(|z_2|^2+|z_3|^2+|z_4|^2+|z_5|^2,
|z_0|^2+|z_1|^2-2|z_5|^2\bigr).
\end{aligned}
\]
$A_3^+(X_\Sigma)\otimes\mathbb R\subset A_3(X_\Sigma)\otimes\mathbb R\simeq\mathbb R^2$ 的 GKZ 分解特别简单，因为生成元包含在 $\Sigma(1)$ 中的射影单纯扇只有两个：$X$ 的扇 $\Sigma$，以及加权射影空间 $\mathbb P(1,1,2,2,2)$ 原来的扇 $\Sigma'$。留给读者作为练习，验证 $A_3^+(X_\Sigma)\otimes\mathbb R$ 的 GKZ 分解如下：
\begin{figure}[H]
\refstepcounter{equation}\label{eq:weighted-gkz-picture}
\centering
\makebox[\textwidth]{\makebox[0pt][l]{(\theequation)}\hfill
\raisebox{-.5\height}{\includegraphics[width=0.56067\textwidth]{3-1-gkz-decomposition.png}}\hfill\mbox{}}
\end{figure}

转向次级扇，注意 $\Xi=\{v_0,\ldots,v_5\}$，因此由式~\eqref{eq:extended-rays}，$\Xi^+$ 由七个点 $(v_0,1),\ldots,(v_5,1),(0,1)$ 组成。它们之间的所有关系构成的空间有一组基，
\SourcePage{74}
由以下矩阵的各行给出：
\[
\begin{pmatrix}
0&0&1&1&1&1&-4\\
1&1&0&0&0&-2&0
\end{pmatrix}.
\]
由\hyperref[subsec:secondary-fan]{第 3.4.2 小节}，在 $A(\Xi)\simeq\mathbb R^2$ 的对偶基下，此矩阵的各列给出次级扇各边的生成元。事实上，从式~\eqref{eq:weighted-gkz-picture} 可见，上述矩阵的前六列给出 GKZ 分解，而最后一列通过沿 $r_1$ 轴负半轴插入一条边，将其扩充为完备扇。
\end{example}

\begin{example}\label{ex:cube-octahedron}
设 $M=\mathbb Z^3$，考虑以原点为中心、顶点为 $(\pm1,\pm1,\pm1)$ 的立方体 $\Delta\subset M_{\mathbb R}$。它给出环面簇 $X=\mathbb P_\Delta$。为描述 $X$ 的扇，注意极多胞形 $\Delta^\circ\subset N_{\mathbb R}$ 是以 $\pm e_1,\pm e_2,\pm e_3$ 为顶点的八面体。因此，法扇由八面体的各面构成，得到扇 $\Sigma$，其三维锥就是 $\mathbb R^3$ 的各卦限。容易推出 $X=\mathbb P^1\times\mathbb P^1\times\mathbb P^1$。此外，来自 $\Delta$ 的除子 $\mathcal O_X(1)$ 是反典范除子 $-K_X$。特别地，$X$ 是 Fano 簇。

$\Delta$ 与 $\Delta^\circ$ 的图形如下：
\begin{figure}[H]
\centering
\includegraphics[width=0.75174\textwidth]{3-2-cube-octahedron.png}
\end{figure}

容易验证，立方体 $\Delta\subset M_{\mathbb R}$ 与八面体 $\Delta^\circ\subset N_{\mathbb R}$ 是\hyperref[sec:fano-reflexive]{第 3.5 节}意义下的一对对偶\IndexLocation{indexloc:1222:74}自反多胞形。特别地，$\Delta^\circ$ 给出“对偶”环面簇 $X^\circ=\mathbb P_{\Delta^\circ}$，它由 $\Delta^\circ$ 的法扇确定；该法扇就是在 $M_{\mathbb R}$ 中向立方体 $\Delta$ 的各面作锥所得的扇。因此得到一对“对偶”环面簇 $X$ 与 $X^\circ$。有趣的是，$X$ 光滑，而 $X^\circ$ 却有相当严重的奇点。事实上，$\Sigma^\circ$ 的三维锥甚至不是单纯的：它们都是无限延伸的四棱锥。不过，$\Delta^\circ$ 自反，所以知道 $X^\circ$ 是 Fano 簇，而且 $\Delta^\circ\subset N_{\mathbb R}$ 是反典范除子 $-K_{X^\circ}$ 对应的多胞形。

还应注意，$\Delta$ 与 $\Delta^\circ$ 的格点数及格点几何分布也很不相同。对 $\Delta^\circ\subset N_{\mathbb R}$，属于 $N$ 的格点只有原点与各顶点；而 $\Delta\subset M_{\mathbb R}$ 中的格点要多得多，因为各棱的中点与各面的中心也属于 $M$。

\SourcePage{75}
还可以描述 $X^\circ$ 的 GKZ 分解。分解的极大锥对应于一维锥生成元位于 $\Sigma^\circ(1)=\{\text{立方体的顶点}\}$ 中的射影单纯扇。对原扇作单纯细分时，需要沿一条对角线将立方体的每个面分成两个三角形。每个面有两种分法，因此 GKZ 分解有 $2^6=64$ 个锥。不过，\hyperref[chap:mirror-constructions]{第 4 章}将说明，为了研究镜像对称，需要使用一维锥生成元位于 $\Delta\cap M$ 中的射影单纯扇来形成 GKZ 分解，而 $\Delta\cap M$ 远大于仅由立方体顶点组成的集合。因此，所得 GKZ 分解的锥数远多于 $2^6$。

这个例题表明，可以从一个光滑环面簇，例如 $\mathbb P^1\times\mathbb P^1\times\mathbb P^1$，出发，通过“对偶”过程得到有奇点的对象。\hyperref[chap:quintic]{第 2 章}中五次三维簇及其镜像出现类似情形，绝非偶然。\hyperref[chap:mirror-constructions]{第 4 章}将详细研究这一现象。
\end{example}
